% ============================================================
%  R. Nielsen, "Om Theologiens Naturbegreb, med særligt Hensyn til
%  Malebranche: De la recherche de la vérité"
%  Kjøbenhavn: Schultziske Officin, 1855
%  Indbydelsesskrift til Kjøbenhavns Universitets Aarsfest
%  til Erindring om Kirkens Reformation
%  TRANSCRIPTION (Danish) — from Google Books scan
%
%  Scan: ~/bibliotek/Nielsen, Rasmus/theologiens_naturbegreb.pdf (58 PDF pp.)
%  TYPE: Fraktur (Danish body); French/Latin quotations in roman/antiqua.
%        OCR with the tessdata_best `Fraktur` script model (NOT -l dan),
%        then image-verify every page (letterspacing -> \emph{}).
%  OFFSET (VERIFIED): PDF = printed + 6.
%        Title (printed p.1)  = PDF 7;  body text begins printed p.3 = PDF 9;
%        last text page printed p.47 = PDF 53 (closing invitation + Univ. seal).
%  EXTENT: continuous essay, printed pp. 3--47 (~45 pp.), no chapter divisions
%        seen; ends with the Indbydelsesskrift's formal invitation + seal.
% ============================================================
\documentclass[12pt,a4paper]{article}
\usepackage[utf8]{inputenc}
\usepackage[T1]{fontenc}
\usepackage{libertinus}
\usepackage{libertinust1math}
\usepackage{textalpha}   % polytonic Greek in text (p.34 Aristotle quote, ὁμοούσιος, θεία φύσις)
\usepackage[danish]{babel}
\usepackage[protrusion=true,expansion=true]{microtype}
\usepackage{setspace}
\usepackage[margin=1.2in]{geometry}
\usepackage{fancyhdr}
\usepackage{hyperref}

\hypersetup{
  pdftitle={Om Theologiens Naturbegreb},
  pdfauthor={R. Nielsen},
  hidelinks
}

\pagestyle{fancy}
\fancyhf{}
\rhead{Nielsen, \emph{Om Theologiens Naturbegreb} (1855)}
\cfoot{\thepage}

\setstretch{1.3}

\title{\textbf{Om Theologiens Naturbegreb}\\[4pt]
       \large med særligt Hensyn til\\[4pt]
       \large Malebranche: \emph{De la recherche de la vérité}}
\author{R.\ Nielsen}
\date{I Anledning af Reformationsfesten 1855.\\[4pt]
  \small Kjøbenhavn: Schultziske Officin, 1855}

\begin{document}
\maketitle
\thispagestyle{fancy}

\bigskip

% === TEXT BEGINS HERE (printed p. 3 = PDF 9) ===

% p. 3
\noindent
Imellem Theologien og Naturlæren er der ingen Forstaaelse; Theologiens
Naturbegreb er ikke naturligt, og Naturlærens ikke theologisk.

Ved endog blot at indskrænke sig til en faktisk Sammenstilling af de
Naturbegivenheder, i hvis Udtydning Naturlæren paa det Bestemteste
modsiger Theologien, vilde det være let at gjøre Uenigheden mellem begge
Videnskaber iøinefaldende. Man røre hvilket Dogme, hvilken theologisk
Forudsætning, man vil; forsaavidt der i nogen Maade kan blive Spørgsmaal om
et virkeligt, tilsvarende Naturforhold, maa Naturlæren i Kraft af
Kjendsgjerningernes Undersøgelse afvise den theologiske Forklaring.

Er det Skabelsen i 6 Dage: Naturlæren beviser, at der alene til Jordens
Dannelse maa være medgaaet Tidsrum af mere end $6 \times 6$ Millioner Aar.

Er det Dogmet om Paradiis og Syndefald: Naturlæren beviser, at der var
Forkrænkelighed, Smerte og Død i Verden, længe før de første Mennesker
bleve til.

Er Talen om Syndfloden og Noæ Ark: Naturlæren beviser, at „Syndfloden" er
en naturlig Local-Catastrophe, og at Dyrene umulig kunne være reddede i Noæ
Ark, eftersom Rovdyrene i faa Dage maatte have fortæret de andre Dyr, og
derefter været døde af Hunger.

Saaledes kunde man blive ved, og endda ikke komme til Maalet. Selv om
Sammenstillingen blev nok saa gjennemført, Grunde og Modgrunde fra begge
Sider fremsøgte og afveiede nok saa omhyggelig; vilde det dog ad denne Vei
neppe være muligt at komme til noget afgjørende Resultat.

% p. 4
Vistnok er Naturlæren med Hensyn paa Sagens Hovedpunkter i det Hele enig med
sig selv, ligesom ogsaa Theologien i det Hele maa siges at følge en til dens
eiendommelige Bestræbelser svarende Retning; men derfor tør man dog ikke
mene, at enhver af de to Videnskaber skulde til en given Tid være ganske og
aldeles enig med sig selv i Udtydning af det Enkelte.

Naturvidenskaben er, som det hedder, i Vorden; denne Vorden, hvorunder den,
hvad Enkelthederne angaaer, stadig corrigerer sig selv, sikkrer sin Methode,
renser og udvider sit Grundlag, er netop betegnende for dens frodige Væxt,
dens rastløse, imponerende Fremskridt.

Theologien er da paa sin Vis ogsaa i Vorden, en Vorden nemlig, hvorunder
dens stedse tiltagende Principløshed og dens Lemmers indvortes Splidagtighed
saa noget nær tyder hen paa Alderdoms Svaghed og tilstundende
Opløsning\footnote{For at det ikke skal synes en eller anden, med de herunder
hørende Forhandlinger mindre bekjendt, Læser, som om denne Paastand
henkastedes uden al Begrundelse, maa jeg indtil videre henvise til mit
Skrift: \emph{Evangelitroen og Theologien}. Jeg kan med saa meget desto
større Føie gjøre det, som ingen Theolog endnu, saavidt mig er bekjendt, har
fundet sig foranlediget til at vise det Ugrundede i den i benævnte Skrift
fremsatte Anskuelse.}.

Ere begge Videnskaber nu saaledes hver i sin Vorden, da skulde det vist kun
lidet hjælpe, om Nogen forfra begyndte enten at veie Grundene for og imod
„Syndfloden", eller at drøfte Problemet om „Noæ Ark", eller, med andre
Ord, at stille Naturlærens Opfatning af enkelte Phænomener imod Theologiens
enkelte Dogmer.

Skal en Forstaaelse mellem Theologi og Naturlære tænkes mulig, maa den vel
begynde med, at de to Videnskaber blive enige i Bestemmelsen af det Begreb,
som for begges videnskabelige Realitet er et afgjørende Grundbegreb, nemlig
Begrebet: \emph{Natur}.

Ved en saadan Begrebsbestemmelse kan her nu ikke nærmest tænkes paa en eller
anden Definition i Almindelighed; thi med formelle Definitioner ud\-

% p. 5
rettes kun lidt i Naturvidenskaben. Opgaven er, at angive det Synspunkt,
hvorunder Naturens Væsen kan sees i det reneste, sandeste videnskabelige Lys.

Stolende paa Sandsernes Vidnesbyrd og Tankens Gyldighed, begynder
Naturlæren da med at tage de tilværende Ting, som de ere givne, skrider saa
ved nøiagtig Prøvelse frem fra det Bekjendte til det endnu Ubekjendte, og
indskrænker sig overalt til at see Naturen i Naturens Lys.

Theologien vælger sig derimod et høiere Udgangspunkt. Begyndende
med Skaberen, gaaer den gjennem Skabelsen over til Skabningen, og kommer
saa, efter egen Formening, til at skue Naturen som Guds Værk i Guds Lys.

Dersom det nu virkelig forholder sig saaledes, at der i Sandhed gives
en Videnskab, som formaaer at fremstille Naturen i Guds Lys, og dersom
denne Videnskab er Theologien: ja, da er det theologiske Naturbegreb, hvad der
saa end kan siges for eller imod en Udtydning af et enkelt Phænomen, unegtelig
den høieste Norm, hvorefter Erfaringsvidenskaben i sin Bedømmelse af
Naturens Væsen sig haver at rette.

Men skulde det derimod lade sig videnskabelig godtgjøre, at der hverken
er eller kan gives en Videnskab, der paa nogen Maade i nogen Grad
formaaer at opklare det Naturlige ved Hjælp af det Overnaturlige: turde vi da
ikke heraf slutte, at den theologiske Skabelseslære med dens tilsvarende
Naturforklaring helst burde anvises en egen Plads udenfor Videnskaben?

Som Repræsentant for den theologiske Opfatning af Naturens Væsen
vælge vi til Sagens Opklaring en skarpsindig Tænker, Forfatteren af Skriftet:
\emph{De la Recherche de la Vérité}\footnote{Paris MDCLXXVIII.}, den ædle
Malebranche.

At Malebranche er Catholik, og hans cartesianske Naturphilosophi længst
antiqveret, faaer i denne Sammenhæng ingen Betydning.

Her er, hvad Dogmerne angaaer, kun eet Dogma, hvorpaa det væsent%
% p. 6
lig kommer an; det er Dogmet om Skabelsen. Her er, hvad Naturlæren
angaaer, kun een Opgave, hvorpaa Opmærksomheden fæstes; det er den
Opgave, at erkjende Naturen som Guds Værk, at skue alle Ting i Gud.

Inderlig forvisset om, at denne store Opgave lader sig løse, er M.
saa langt fra at trække sig tilbage i det Mysteriøse, og skjule Tankernes
Fattigdom i æsthetisk-taagede Speculationer, at han meget mere opfordrer til at
begynde med det Enkelte, det Tydelige, det let Forstaaelige, og fører paa
mange Steder et Sprog, saa man skulde troe, at det var, ikke En af
Oratoriets Fædre, der talede om Guds Almagt, men En af Nutidens
Naturkyndige, der afviste philosophisk Mysticisme\footnote{Jfr. t.~Ex. De l.
Recherch. S.~394. Combien de gens rejettetent la Philosophie de M. Descartes
par cette plaisante raison que les principes en sont trop simples et trop
faciles. Il n'y a point de termes obscurs et mysterieux dans cette Philosophie:
des femmes et des personnes qui ne sçavent ni grec ni latin, sont capables de
l'apprendre: il faut donc que soit peu de chose, et il n'est pas juste que de
grands genies s'y appliquent. Ils s'imaginent que des principes si clairs et si
simples ne sont pas assez féconds, pour expliquer les effets de la nature qu'ils
supposent obscure et embarassée. \dots\ Ils aiment donc mieux expliquer des
effets, dont ils ne comprennent point la cause, par des principes qu'ils ne
conçoivent point, et qu'il est absolument impossible de concevoir, que par des
principes simples et intelligibles tout-ensemble. Car ces Philosophes expliquent
des choses obscures par des principes qui ne sont pas seulement obscurs, mais
entièrement incompréhensibles.}.

\begin{center}---\end{center}

„Det er" --- saa tænker Malebranche --- „Vildfarelsen, der er Aarsag til
Menneskets Elendighed; det er derfor nødvendigt, at den af alle Kræfter
modarbeides. Man maa ikke indbilde sig, at det er saa besværligt at søge
Sandheden. Man behøver kun at aabne sine Øine, være opmærksom, og følge
nogle Regler, som vi i det Følgende ville give. Det er jo dog den skarpe
Tænkning, der bringer Lys, og opdager Sandheden".

„Endskjøndt vore Sandser ere fordærvede ved Synden, er det dog
egentlig ikke Sandserne, men vor Frihed, der er Aarsag i vore Vildfarelser".

% p. 7
„Da vi nemlig ere sammensatte af Aand og Legeme (Etant composez
d'un esprit et d'un corps), have vi to Slags Goder at søge, Aandens og
Legemets. Vi have da ogsaa to Maader, hvorpaa vi kunne kjende, om Noget
er godt eller ondt. Vi kunne kjende et Gode enten ved en klar Indsigt eller
en uklar Sandsning (un sentiment confus). Jeg erkjender med Fornuften,
at Retfærdighed er elskelig; jeg veed ogsaa af Smagen, at denne eller hiin
Frugt er god. Retfærdighedens Skjønhed kan ikke sandses (ne se sent pas),
Frugtens Godhed kan ikke tænkes (ne se connoit pas). Legemets Goder
fortjene da heller ikke, at en Aand, som Gud kun har skabt for sig, randsager
dem; det er derfor nødvendigt, at Aanden uden nærmere Undersøgelse (examen)
maa kunne kjende slige Goder ved Sandsningens korte, utvivlsomme
Prøve\footnote{S.~18. Les pierres ne sont pas propres à la nourriture, la preuve
en est convaincante, et le seul goût en a fait tomber d'accord tous les
hommes.}".

„Dersom Aanden kun saae i Legemerne, hvad der virkelig er i dem,
saa kunde den ikke elske dem og ikke uden megen Uvillie betjene sig af dem;
det er da nødvendigt, at de synes behagelige, idet de foraarsage Følelser, der
ikke vedkomme deres Væsen".

„Men saaledes er det ikke med Gud, som alene er Aandens sande
Gode; ham er det ikke nok at elskes med Instinctets blinde Kjærlighed; han
vil elskes med en oplyst, klar og besluttet Kjærlighed (d'un amour éclairé et
d'un amour déchoix)".

„For ret at forstaae dette, maa man vide, at der er en meget betydelig
Forskjel mellem den Maade, hvorpaa Naturens Ophav (l'Auteur de la Nature)
bevæger Materien, og den Maade, hvorpaa han uden Ophør bevæger Aanden
mod det Gode i Almindelighed (vers le bien en general). Thi Materien
er uden al Virksomhed, uden Kraft til at standse sin Bevægelse eller bestemme
dens Retning."

„Anderledes med Villien; om den kan man dog i en vis Forstand (en
un sens) sige, at den er virksom, og at den i sig selv har Kraft til at
% p. 8
optage det Indtryk (impression), Gud giver den\footnote{S.~5. car quoi qu' elle
ne puisse pas arrêter cette impression, elle peut en un sens la détourner du
côté qu'il lui plaît\dots}, paa forskjellig Maade,
og derved foraarsage al den Uorden (déreglement), der viser sig i
Tilbøielighederne, og alle de Elendigheder, der ere Følger af Synden".

„Er end Villien ikke fri i den Forstand, at den skulde have Magt til
at ville og ikke at ville, saa er den dog istand til at drive os frem mod det
Gode i Almindelighed\footnote{S.~5. \dots\ par ce mot de \emph{Volonté}, je
prétens ici désigner l'impression ou le mouvement naturel, qui nous porte vers
le bien indéterminé et en général: et par celui de \emph{Liberté}, je n'entens
autre chose que la force qu' a l'esprit de détourner cette impression vers les
objets, qui nous plaisent.}".

„Da Villien altid retter sig efter Forestillingen om det Gode, og vi
umulig kunne ville Det, vi i samme Øieblik maae ansee for et væsentligt
Onde, saa maa der være den inderligste Forbindelse mellem Erkjendelse og
Villie, mellem det Sande og det Gode\footnote{S.~6. C'est à-peu-près la même
chose de la connoissance de la vérité que de l'amour du bien.}".

„Sandser og Lidenskaber (les sens et les passions) have ikke deres
Oprindelse fra Synden, men den Magt, hvormed de tyrannisere Syndere.
Denne Magt er ikke saameget en Forstyrrelse paa Sandsernes Side, som en
Forvirring i Menneskets Aand og Villie".

„Adam havde i sin Uskyldighedsstand samme Sandser som vi, men
var ubetinget Herre over de Indtryk af Glæde og Smerte, de opvakte i hans
Legeme. Først, da han frivillig vendte sig bort fra Guds Aasyn, for ligesom
at mætte sin Aand med den forbudne Frugts Skjønhed og forventede Sødme,
gjorde hans Sandser og Lidenskaber Oprør imod ham, og gjorde ham, som
os, til en Slave af alle sandselige Ting".

„Vore Sandser ere ikke saa fordærvede, som man indbilder sig; men
det er vor Sjæls Inderste, det er vor Frihed, der er fordærvet. Det er
% p. 9
ikke vore Sandser, der bedrage os; det er vor Villie, der bedrager os ved
dens overilede Domme".

„Naar man t.~Ex. seer Lys, da er det ganske vist, at man seer Lys;
naar man føler Varme, da bedrager man sig ikke ved at troe, at man føler
den --- være sig før eller efter Syndefaldet".

„Men man bedrager sig, naar man antager (quand on juge), at den
Varme, man føler, er udenfor Sjælen, som føler den. Det er en overilet
Dom, der beroer paa Frihedens Misbrug. Derfor bør vi nøie iagttage den
Regel, aldrig at dømme efter Sandserne om, hvad Tingene ere i sig selv;
thi Sandserne ere i Virkeligheden ikke givne, for at vi skulle kjende Tingenes
Sandhed i sig, men alene for at bevare vort Legeme (pour la conservation
de notre corps)".

„Vildfarelsens Grund ligger heller ikke i vor Forstands Indskrænkning;
thi vel er Forstanden indskrænket, men det er dog Villien, der leder vor
Dømmekraft\footnote{S.~6 \dots\ nous sommes aussi libres dans nos faux
jugemens, que dans nos amours déreglez\dots}. Det bliver derfor en Villiens
Opgave, at arbeide sig ud af Blændværket og søge Sandheden. Men at søge
Sandheden er i alle Ting at søge Gud".

„Sandheden kjendes paa Klarheden, paa det Indlysende\footnote{S.~9. Car
la vérité ne se trouve presque jamais qu' avec l'évidence, et l'évidence ne
consiste que dans la vûe claire et distincte de toutes les parties et de tous les
rapports de l'objet, qui sont nécessaires pour porter un jugement assûré.}; det
maa derfor indskjærpes, at vi aldrig indrømme Noget, før det er os saa
indlysende, at vi tvinges ved vor Fornufts indre alvorlige Tilretteviisning
(jusqu'à ce que nous y soyons comme forcez par des reproches intérieures de
notre raison\footnote{S.~1--2. L'exactitude de l'esprit n'a presque rien de
penible: ce n'est point une servitude comme l'imagination la représente; et si
nous y trouvons d'abord quelque difficulté, nous en recevons bientôt des
satisfactions qui nous récompensent bien de nos peines; car enfin il n'y a
qu'elle qui produise la lumiére, et qui découvre la vérité.})".

% p. 10
Efter saaledes at have paaviist Vildfarelsens Aarsag, og derhos (1ste
og 2den Bog) udførlig afhandlet de ved Sandserne og Indbildningskraften
foranledigede Vildfarelser, gaaer Forfatteren (3die Bog) over til at omhandle
Aandens Væsen: \emph{De l'entendement ou de l'esprit pur}.

Det almindelige Forhold mellem Aand og Materie er allerede i første
Bog opfattet paa følgende Maade:

„Ligesom Materien, det Udstrakte, har et dobbelt Anlæg, Anlæg til at
modtage forskjellig Form, og Anlæg til at bevæges; saaledes har ogsaa
Menneskeaanden to Evner, \emph{Forstanden} (l'entendement), hvormed vi
opfatte forskjellige Ideer, det vil sige: bringe forskjellige Ting til Bevidsthed,
og \emph{Villien} (la volonté), Anlæget for mangfoldige Tilbøieligheder eller
det, at ville forskjellige Ting".

Denne Sammenstilling af det Materielle og det Aandelige, hvorved
tilsyneladende Lighedspunkter fremhæves, maa dog ikke forstaaes, som om her
forudsattes et naturligt Slægtskab mellem Aand og Materie. Tvertimod!
det indskjærpes gjennem hele Undersøgelsen fra Først til Sidst, at de to
Substanser ere saa ueensartede, som muligt.

Uagtet der tillægges dem forskjellige Evner, have de dog alligevel
ingen virksomme Evner\footnote{S.~4. C'est pour nous accommoder à la maniére
ordinaire de parler, que nous dirons dans la suite, que les sens sentent\dots}.
Begrebet: „Evne" maa her saavel for Aandens som Materiens Vedkommende tages
i en saa passiv Forstand, at det slet ikke udtrykker Anlæg til gjensidig
Virksomhed, men kun en Mulighed til at paavirkes af Gud\footnote{S.~5. De même
que l'Auteur de la Nature est la cause universelle de tous les mouvemens, qui se
trouvent dans la matiére; c'est aussi lui qui est la cause générale de toutes les
inclinations naturelles qui se trouvent dans les esprits\dots}.

„Aandens Væsen", hedder det nu (3die Bog), „bestaaer alene i
Tænkning, ligesom Materiens alene i Udstrækning. Det er ifølge Tænkningens
% p. 11
forskjellige Modificationer, at Aanden snart er villende, snart forestillende,
eller har flere andre særegne Former".

„Det beroer ligeledes paa forskjellige Modificationer, nemlig af
Udstrækning, at Materien er snart Vand, snart Træ, snart Ild, eller har en
Uendelighed af andre saadanne Former\footnote{S.~171. Je ne croi pas aussi
qu'il soit possible de concevoir un esprit qui ne pense point, quoi qu'il soit
fort facile d'en concevoir un qui ne sente point, qui n'imagine point, et même
qui ne veüille point: de même qu'il n'est pas possible de concevoir une matiére,
qui ne soit pas étendue quoi qu'il soit assez facile d'en concevoir une, qui ne
soit ni terre, ni métal ni quarrée ni ronde et qui même ne soit point en
mouvement. Il faut conclure de là, que comme il se peut faire qu'il y ait de la
matiére, qui ne soit ni terre ni metal, ni quarrée ni ronde ni même en
mouvement: il se peut faire aussi qu'un esprit ne sente ni chaud ni froid, ni
joie, ni tristesse, n'imagine rien, et même ne veuille rien; de sorte que toutes
ces modifications ne lui sont point essentielles. \emph{La pensée toute seule
est donc l'essence de l'esprit, ainsi que l'étendue toute seule est l'essence de
la matiére.}}".

Da nu Aandens Væsen sættes i Tanken, kunde det jo vistnok synes, som om her
i det Hele blev lagt et for stærkt Eftertryk paa Villien; men M. bemærker
udtrykkelig, at Evnen til at ville alligevel er uadskillelig fra Aanden,
endskjøndt den ikke udgjør dens Væsen, ligesom Evnen til at bevæges er
uadskillelig fra Materien, uagtet den heller ikke udgjør Materiens Væsen.

„Er end Tænkningen Aandens Væsen, tør vi dog ingenlunde deraf slutte, at den
tænkende Aand skulde være istand til selv at kjende alle de Modificationer, hvis
Mulighed den indeholder\footnote{S.~173. Un simple morceau de cire est donc
capable d'un nombre infini, ou plûtôt d'un nombre infiniment infini de
différentes modifications, que nul esprit ne peut comprendre. Quelle raison donc
de s'imaginer que l'âme qui est beaucoup plus noble que le corps, ne soit
capable que des seules modifications qu'elle a déja reçuës.}".

„Naar vor Sjæl hernede kun modtager meget faa Modificationer, da er det, fordi
den er forbunden med et Legeme, og afhængig deraf".

„Den Materie, hvoraf vort Legeme er sammensat, er kun istand til
% p. 12
at modtage meget faa Modificationer, saalænge vi leve. Denne Materie kan ikke
opløse sig til Jord, til Damp, før efter vor Død. Den kan i vor Levetid ikke
blive til Luft, Ild, Diamant, Metal; den kan ikke blive rund, fiirkantet,
trekantet: den maa være Kjød, og have en menneskelig Figur, for at vor Sjæl
kan være forenet med den".

„Paa samme Maade forholder det sig med Sjælen. Det er nødvendigt, at den har
Sandsninger af Varme, af Kulde, af Farve, af Lys, af Toner, af Lugt, af Smag og
mange andre Modificationer. Disse Sandsninger lære den just at bevare sin
Maskine (l'appliquent à la conservation de sa machine\footnote{Elles l'agitent,
et l'effraient dés que le moindre ressort se débande, ou se rompt, et ainsi il
faut que l'âme y soit sujette, tant que son corps sera sujet à la corruption.
Mais lorsqu'il sera revêtu de l'immortalité, et que nous ne craindrons plus la
dissolution de ses parties, il est raisonable de croire qu'elle ne sera plus
touchée de ces sensations incommodes que nous sentons malgré nous; mais d'une
infinité d'autres toutes différentes, dont nous n'avons maintenant aucune
idée.})".

„Man vogte sig vel for at adsprede Aanden ved at ville forklare Ting, der ere
altfor indviklede. Menneskeaanden er nemlig underkastet Vildfarelser, ikke blot
fordi den ikke er uendelig, men ogsaa fordi den ikke har Udholdenhed nok til at
dvæle ved hver enkelt Gjenstand saalænge, indtil denne i alle Henseender er
fuldstændig og nøiagtig undersøgt".

„De Gjenstande, hvormed Aanden beskjæftiger sig, kunne inddeles i to Classer:
enten ere de i Sjælen eller udenfor den".

„De, der ere i Sjælen, ere dens egne Tanker, hvorved i Almindelighed forstaaes
alle de Ting (toutes les choses!), der ikke kunne være i Sjælen, uden at den
bliver dem vaer, saasom: dens egne Sandsninger, dens Imaginationer, rene
Begreber og Forestillinger, selv dens Lidenskaber (ses passions mêmes) og
naturlige Tilbøieligheder".

Hvad disse Gjenstande angaaer, kan Sjælen vel erkjende dem uden at behøve et
Mellemværende; thi de ere jo ikke Andet end Sjælen selv, modi%
% p. 13
ficeret paa denne eller hiin Maade (d'une telle ou telle façon); men
\emph{hvorledes kommer Sjælen til Erkjendelse af de Gjenstande, der ere udenfor
Sjælen?}

Det er et Hovedspørgsmaal.

„Vi see Solen, Stjernerne, en uendelig Mængde Gjenstande udenfor os; det er dog
ikke sandsynligt, at Sjælen gaaer ud af Legemet, at den, saa at sige, gjør sig
en Spadseretour gjennem Himlene, for der at beskue alle disse Gjenstande. Den
seer dem da heller ikke ved dem selv. Aandens umiddelbare Object, naar den seer
Solen, er ikke Solen selv, men et Noget, der er inderlig forenet med vor Sjæl,
og det er det, jeg kalder Ideen" (S.~188. \emph{De la nature des idées})".

Det Charakteristiske ved M.'s Ideelære er nu en ligefrem Følge af den
eiendommelige Maade, hvorpaa han, i Forhold til Datidens philosophiske
Standpunkt, forklarer sig Erkjendelsen af de udvortes Gjenstande.

„\emph{Materielle Gjenstande kunne umulig opfattes ved dem selv; de opfattes kun
ved de i Sjælen tilsvarende Ideer: hvorfra komme da disse Ideer?}\footnote{S.~190.
Nous assurons donc qu'il est absolument nécessaire, que les idées, que nous avons
des corps, et de tous les autres objets, que nous n'appercevons point par eux
mêmes, viennent de ces mêmes corps ou de ces objets: ou bien que nôtre âme ait
la puissance de produire ces idées: ou que Dieu les ait produites avec elle en
la créant, ou qu'il les produise toutes les fois qu'on pense à quelque objet: ou
que l'âme ait en elle-même toutes les perfections qu'elle voit dans ces corps:
ou enfin qu'elle soit unie avec un être tout parfait, et qui renferme
généralement toutes les perfections des êtres créez.}".

Til en Besvarelse af dette Spørgsmaal baner M. sig Veien gjennem en Kritik af de
fra hans Anskuelse afvigende Theorier.

„Peripatetikernes Mening, at de udvortes Gjenstande skulde tilsende os deres
Afbilleder (des espèces, qui leur ressemblent), og at disse saa igjennem de ydre
Sandser skulde bringes til den fælleds Sands (au sens commun),
% p. 14
er urimelig; disse Indtryksbilleder (especes impresses) maatte da selv være smaa
Legemer\footnote{S.~191. Tous les objets, comme le Soleil, les Etoiles, et tous
ceux qui sont proche de nos yeux, ne peuvent pas envoyer des especes, qui soient
d'autre nature qu'eux \dots\ c'est pourquoi les Philosophes disent ordinairement,
que ces especes sont grossieres et matérielles, à la différence des especes
\emph{expresses}: (especes matérielles et sensibles, qui sont rendues
intelligibles par l'intellect agent ou agissant, et sont propres pour être
reçues dans l'intellect patient S.~190) qui sont spiritualisées.}, og disse
Legemer, hvoraf Rummet vilde være opfyldt, maatte da ifølge Materiens
Uigjennemtrængelighed knuse og støde hverandre, idet nogle gik til en, andre til
en anden Side, og kunne altsaa ikke gjøre Gjenstandene synlige".

Deres Mening, som antage, at Gjenstandenes Ideer ere skabte med os, saa at Gud
ikke behøver i hvert Øieblik at vække dem i Sjælen, naar vi have dem behov,
gjendrives ved følgende Betragtning:

„Vi have, selv hvad simple Figurer angaaer, en uendelig Mængde Ideer. Ellipser,
Triangler, Polygoner o. desl. kunne i Henseende til Størrelsesforhold imagineres
paa utallige Maader\footnote{S.~196. L'esprit voit donc toutes ces choses: il en
a des idées. Il est sûr que ces idées ne lui manqueront jamais, quand il
employeroit des siécles infinis à la considération même d'une seule figure.}. Det
er da høist usandsynligt, at Gud skulde have skabt en saadan Vrimmel af Ideer i
og med den menneskelige Aand. Men sæt, at vi ogsaa havde et saadant Oplag af
medskabte Ideer (un magazin de toutes les idées), saa bliver det dog aldeles
uforklarligt, hvorledes vi komme til at anvende den rette Idee, naar den
tilsvarende Gjenstand viser sig. Vi see t.~Ex. Solen. Det Billede, Solen præger
i Hjernen, har, som allerede viist, ingen Lighed med den Idee, vi danne os om
Solen; heller ikke mærker Sjælen det allermindste til det Indtryk, Solen gjør
enten paa Øiets Grund eller i Hjernen: det er altsaa umuligt, at forklare,
hvorledes vi blandt de utallige Ideer just komme til at vælge Solens Idee, naar
vi see Solen".
% p. 15
„De, som paastaae, at vore Sjæle selv have den Evne at kunne danne sig Ideer
svarende til de Gjenstande, hvorover vi ville tænke, skjøndt de fra disse
Gjenstande kommende Indtryk ikke ere Billeder, der ligne Gjenstandene, have
enten altfor ringe Forestillinger om Ideerne eller altfor høie Tanker om den
menneskelige Aand\footnote{S.~193. En effet le monde intelligible doit être plus
parfait que le monde matériel et terrestre\dots\ Ainsi, quand on assure, que les
hommes ont la puissance de se former les idées telles qu'il leur plaît, on se met
fort en danger d'assurer que les hommes ont la puissance de faire des êtres plus
nobles et plus parfaits que le monde que Dieu a créé. On ne fait pas cependant
réflexion à cela, parce qu'on s'imagine, qu'une idée n'est rien, à cause qu'elle
ne se fait point sentir \dots}".

„En saadan Indbildning om Sjælens Fuldkommenhed er formastelig. Kun Gud alene
kan i sig selv skue og erkjende alle Ting, efterdi han er alle Tings Skaber.
„Siger ikke, at I have Lyset i Eder selv, siger den hellige Augustinus, thi kun
Gud alene har Lyset i sig selv".

„Saa staaer der da kun een Anskuelse tilbage, som den eneste fornuftige (qui
paroit seule conforme à la raison), den nemlig, at Aanderne have alle deres
Tanker fra Gud".

„For ret at fatte dette, maae vi for det Første erindre, at Gud nødvendigviis
maa have alle skabte Tings Ideer i sig selv, da han jo ellers ikke kunde
frembringe dem; dernæst, at Gud ifølge sin Allestedsnærværelse er saa inderlig
forenet med vore Sjæle, at man endog kunde sige, at han er Aandernes Sted
((qu'il est le lieu des esprits), ligesom man siger om Rummet, at det er
Legemernes Sted".

„Under disse Forudsætninger er det altsaa vist, at Aanden kan see, hvad det er,
der er i Gud, som fremstiller de skabte Ting, da det er saa aandigt, saa
intelligibelt, saa nærværende for Aanden. Aanden kan da i Gud see Guds Værk, saa
sandt Gud vil aabenbare den det".
% p. 16
Saaledes kommer M. da til at opstille den for Theologien saa uundværlige
Grundsætning, at vi nemlig skue alle Ting i Gud:

\begin{center}
\textbf{Que nous voyons toutes choses en Dieu.}
\end{center}

Denne Fundamentalsætning maa imidlertid ikke misforstaaes, som om Meningen var,
at vi indskrænkede Mennesker ved at see Tingene i Gud nu ogsaa kunne gjennemskue
Guds Væsen\footnote{S.~200. Mais il faut bien remarquer, qu'on ne peut pas
conclure que les esprits voyent l'essence de Dieu, de ce qu'ils voyent toutes
choses en Dieu de cette maniére. Parce que ce qu'ils voyent est très-imparfait,
et que Dieu est très-parfait.}. Langtfra! Hvorledes skulde det Ufuldkomne kunne
fatte det Fuldkomne?

„En Grund mere til at antage, at vi see alle Ting, fordi Gud vil det, er da
ogsaa dette, at vi derved erkjende vor uendelige Afhængighed af Gud; thi af os
selv vide vi Intet. Non sumus sufficientes cogitare aliquid a nobis, tanquam ex
nobis, sed sufficientia nostra a Deo est. 2 Cor. 3, 5".

„Det er Gud selv, der oplyser Philosopherne i de Erkjendelser, som disse
utaknemmelige Mennesker kalde naturlige, uagtet de komme til dem fra Himlen. Deus
enim illis manifestavit. Rom. 1, 19. Det er ham, der er Aandens sande Lys: Pater
luminum. Jac. 1, 17. Det er ham, der lærer Menneskene Videnskab: Qui docet
hominem scientiam. Psalm. 94, 10. Med eet Ord: det er det sande Lys, som oplyser
alle dem, der komme til denne Verden: lux vera, quæ illuminat omnem hominem
venientem in hunc mundum. Joh. 1, 9".

Ingen vil kunne miskjende, at dette er altsammen theologisk, ægte theologisk.

Disse Træk af en i sig uholdbar Erkjendelseslære kunne imidlertid i denne
Sammenhæng kun interessere os, forsaavidt de afgive Forudsætninger for en
tilsvarende Naturbetragtning, og derved gjøre det indlysende, hvad der
% p. 17
ligger i Theologiens Naturbegreb, naar de i samme indeholdte Conseqvenser
fremsættes saa klart, saa alsidig, saa uforbeholdent, som hos Malebranche.

Er Guds, i de tænkende Aander sig fortsættende, Aabenbaring Erkjendelsens eneste
sande Kilde, saa forstaaer det sig vel af sig selv, at Tingenes Væsen alene
beroer paa den guddommelige Villie, at Intet i Naturen har Kraft og Magt til
selv at virke, men at det maa være Gud, Gud alene, der virker Alt i alle Ting.

„Aristoteles\footnote{S.~383. Car les termes ordinaires à ce Philosophe ne
peuvent servir qu'à exprimer confusément aux sens et à l'imagination les
sentimens confus que l'on a des choses sensibles: ou à faire parler d'une
maniére si vague et si indéterminée, que l'on n' exprime rien de distinct.
Presque tous ses ouvrages, mais principalement ses huit Livres de Physique
\dots\ ne sont qu, une pure Logique.} og de Philosopher, der følge hans Banner,
forfeile Sandheden, og vandre i Mørket. Stolende paa Sandseindtrykkene, give de
os kun tomme Ord og forvirrede Forestillinger om Tingene."

„Alle de Kategorier (termes), som kun betegne sandselige Forestillinger, ere
tvetydige, vel at mærke, tvetydige ifølge Vildfarelse og Uvidenhed, og saaledes
igjen Aarsag til en uendelig Mængde Vildfarelser."

„Ordet „Væder" f.~Ex. er tvetydigt; det betyder et drøvtyggende Dyr; det betyder
ogsaa en Constellation, hvori Solen indtræder om Foraaret. Ved et saadant Udtryk
bliver nu rigtignok Ingen saa let bedragen; thi man maatte da indtil det Yderste
være Astrolog for at indbilde sig, at her kunde være nogen virkelig Forbindelse
imellem Væder og Væder. Men hvad Sandselighedskategorierne ellers angaaer,
bliver man let skuffet; thi der er næsten Ingen, der lægger Mærke til, at de
tilhobe ere tvetydige."

„Naar t.~Ex. Philosopherne sige: Ilden er hed; Planten er grøn, Sukkeret sødt,
o. s. v., forstaae de det ligesom Børn og Hverdagsmennesker, som om det var
Ilden, der indeholdt den Hede, de føle, Planten den Farve, de mene at see, og
Sukkeret den Sødme, de smage, naar de spise det; og
% p. 18
saaledes med alle Ting, vi see eller sandse. De sige ikke to Ord uden at sige
noget Falsk: Alt, hvad de sige om denne Materie, er dunkelt og forvirret."

Dette bevises da af forskjellige Grunde:

„For det Første, ere Menneskene høist uenige om de sandselige Indtryk: hvad den
Ene finder sødt, finder den Anden bittert; hvad den Ene finder varmt, finder den
Anden koldt."

„For det Andet, kunne de allerforskjelligste Gjenstande gjøre samme
Sandseindtryk. Gips, Snee, Sukker, Salt o. s. v. give samme Farvefornemmelse:
desuagtet er disse Tings Hvidhed meget forskjellig, naar man ellers vil dømme
efter Andet end den blotte Sandsning. Naar man altsaa siger: Melet er hvidt,
siger man intet bestemt."

„For det Tredie, ere de legemlige Egenskaber, der foraarsage de meest
forskjellige Sandseindtryk, nu og da næsten de samme, og, omvendt, ere de, der
frembringe omtrent samme Indtryk, ofte meget forskjellige. Qvaliteterne: Sødhed
og Bitterhed ere næsten ikke forskjellige, men Fornemmelserne af Sødhed og
Bitterhed ere væsentlig forskjellige. De Bevægelser (mouvemens), der foraarsage
Smerte og Kildren, ere kun forskjellige ved et Mere og Mindre, men Følelserne af
Kildren og Smerte ere væsentlig forskjellige."

„En anden Række af tvetydige Udtryk hentes fra Logiken, ved hvis Almeenbegreber
man let kan forklare alle Ting uden Kundskab til
Tingene.\footnote{S.~386. Aristote est celui qui en a le plus fait usage, tous
ses Livres en sont pleins, et il y en a quelques uns, qui ne sont que pure
Logique. Il propose et résout toutes choses par ces beaux mots de \emph{genre},
d'\emph{espece}, d'\emph{acte}, de \emph{puissance}, de \emph{nature}, de
\emph{forme}, de \emph{facultez}, de \emph{qualitez}, de \emph{cause par soi}, de
\emph{cause par accident}. Ses sectateurs ont de la peine à comprendre que ces
mots ne signifient rien \dots}"

„Man bliver ikke klogere ved at høre Aristotelikerne sige, at Ilden smelter
Metaller, fordi den har Evne til at smelte, eller at En ikke kan fordøie, fordi
hans Mave ikke har Evne o. s. v.\footnote{S.~386. Le feu échauffe, séche, durcit,
et amollit, parce qu'il a la faculté de produire ces effets. Le sené purge par sa
qualité purgative, le pain même nourrit, si on le veut, par sa qualité nutritive
\dots\ Telles ou semblabes maniéres de parler ne sont point fausses, mais c'est
qu'en effet elles ne signifient rien.}"
% p. 19
„Men hvad der er det Allerfarligste ved disse philosophiske Talemaader:
Qvalitet, Evne, Væsenhed o. desl., er dog dette, at man derved hildes i den
Indbildning, at Tingene selv skulde indeholde de sande og oprindelige Aarsager
til de Virkninger, vi see for vore Øine."

„Det er en syndig Vildfarelse; thi, dersom der i Tingene selv ere virkelige
Aarsager, virkelige Evner til selv at virke, saa er der i Tingene noget
Guddommeligt, og vi komme da, ligerviis som Hedningerne, til at forgude
Naturen.\footnote{S.~387. Que si l'on vient ensuitte à considerer attentivement
l'idée que l'on a de cause ou de puissance d'agir, on ne peut douter que cette
idée ne represente quelque chose de divin. Car l'idée d'une puissance souveraine
est l'idée de la souveraine divinité, et l'idée d'une puissance subalterne est
l'idée d'une divinité inférieure, mais d'une véritable divinité, au moins, selon
la pensée des Payens, supposé que ce soit l'idée d'une puissance ou d'une cause
véritable. On admet donc quelque chose de divin dans tous les corps qui nous
environnent, lorsqu'on admet des formes des facultez, des qualitez, des vertus,
et des êtres réels capables de produire certains effets par la force de leur
nature, et l'on entre ainsi insensiblement dans le sentiment des Payens par le
respect que l'on a pour leur Philosophie. Il est vrai que la foi nous redresse,
mais peut-être que l'on peut dire, que si le coeur est chrétien, le fond de
l'esprit est Payen.}"

„Det er vanskeligt at troe, at man hverken bør frygte eller elske de virkelige
Magter, de Væsner, som kunne indvirke paa os, som kunne straffe os med Smerte og
belønne os med Glæde. Da nu sand Tilbedelse bestaaer i Kjærlighed og Frygt, saa
er det endydermere vanskeligt at troe, at man ikke bør tilbede dem."

„Alt, hvad der kan indvirke paa os som sand og virkelig Aarsag, er nødvendigviis
ophøiet over os (est nécessairement au dessus de nous). Naar derfor Gud i den
hellige Skrift beviser Israeliterne, at de bør tilbede ham, det vil sige, at de
bør frygte og elske ham, da ere de vigtigste Grunde, han anfører, hentede fra
hans Magt til at belønne dem og straffe dem.
% p. 20
Han vil, at de skulle ære ham, efterdi det er ham alene, der er den sande Aarsag
til Godt og Ondt; og efterdi der, ifølge Prophetens Ord, intet Ondt skeer i
deres By, uden at det jo er ham (Jehova), der gjør det.

„De naturlige Aarsager ere altsaa ingenlunde virkelige Aarsager til det Onde, de
synes at paaføre os, eftersom det er Gud alene, der virker i dem, ham alene, vi
bør frygte og elske: Soli Deo honor et gloria."

Dersom dette ikke er theologisk, ægte theologisk, hvad er saa theologisk?

„Antage vi", hedder det endvidere, „Philosophernes falske Mening (cette fausse
opinion des Philosophes), at nemlig Legemerne, som omgive os, ere virkelige
Aarsager til de Glæder og Onder, vi føle; da synes Fornuften at retfærdiggjøre en
Religion lig Hedningenes, og at bifalde en almindelig Sædernes Fordærvelse
(déreglement universel des meurs)."

„Tilsiger Fornuften os end ikke ligefrem, at vi derfor strax skulle tilbede
Zwibler og Porreløg, saa tilraader den os dog at yde dem en vis Dyrkelse i
Forhold til det Gavn, de kunne yde os.\footnote{S.~388. Il est vrai que la raison
n'enseigne pas qu'il faille adorer les oignons et les porreaux par exemple, comme
la souveraine divinité, parce qu'ils ne peuvent nous rendre entiérement heureux,
lorsque nous en avons, ou entiérement malheureux lors que nous n'en avons point
\dots\ Mais (S.~389), si l'on ne doit pas rendre un honneur souverain aux
porreaux et aux oignons, on peut toujours leur rendre quelque adoration
particuliére: je veux dire, qu'on peut y penser et les aimer en quelque maniére,
s'il est vrai qu'ils puissent en quelque sorte nous rendre heureux: on doit leur
rendre honneur à proportion du bien qu'ils peuvent faire.}"

Inderlig oprørt over det Falske i en saa elendig Philosophie (de la fausseté de
cette misérable Philosophie), udvikler Malebranche saa den Lære:
\emph{at der kun er een sand Aarsag, efterdi der kun er een sand Gud; at Naturen
eller Tingenes Kraft ikke er Andet end Guds Villie; at de naturlige Ting ikke ere
sande Aarsager, men aleneste Anledningsaarsager (des causes occasionelles).}
% p. 21
Her have vi da det saakaldte Occasionalsystem: systema causarum occasionalium.

Hvorledes denne Theori i Virkeligheden skal forstaaes, bliver anskueliggjort ved
følgende Exempel:

„Naar en Kugle støder sammen med en anden, og sætter denne i Bevægelse, saa
meddeler den dog alligevel Intet; thi den har ikke selv det Indtryk, den
meddeler. Imidlertid bliver en Kugle dog en naturlig Aarsag til den Bevægelse,
den meddeler. En naturlig Aarsag er altsaa ingen sand og virkelig Aarsag, men
derimod en Anledningsaarsag, som bestemmer Naturens Skaber til at virke paa denne
eller hiin Maade i dette eller hiint Tilfælde.\footnote{Jvnf. S.~390 med S.~575.
Quand je vois une boule qui en choque une autre, mes yeux me disent, ou semblent
me dire, qu'elle est véritablement cause du mouvement qu'elle lui imprime: car la
véritable cause qui meut les corps ne paroît pas à mes yeux. Mais quand
j'interroge ma raison, je vois évidemment que les corps ne pouvant se remüer
eux-mêmes, et leur force mouvante n'étant que la volonté de Dieu qui les conserve
successivement en différens endroits, ils ne peuvent communiquer une puissance
qu'ils n'ont pas, et qu'ils ne pourroient pas mêmes communiquer quand elle seroit
en leur disposition. Car il est évident qu' il faut une sagesse, et une sagesse
infinie pour régler la communication des mouvemens avec la justesse, la
proportion, et l'uniformité que nous voyons.} Lovene, hvorefter Bevægelserne
meddeles, ere kun virksomme, forsaavidt Guds Villie virker i dem. Gud har skabt
Verden, fordi han har villet det: dixit et facta sunt. Der er altsaa i den
sandselige Verden ingen Kræfter, ingen Magter, ingen sande Aarsager, og man bør
ikke ved at tillægge Legemerne Former, Evner og reelle Qvaliteter, de ikke have,
bedrage sig selv, og krænke den guddommelige Majestæt."

„Gud behøver ikke Redskaber for at virke (Dieu n'a pas besoin d'instrumens pour
agir); der behøves kun, at han vil det, for at Tingen kan være. Hans Magt er
hans Villie. Men at meddele sin Magt til et Menneske eller til en Engel, kan
ikke betyde Andet end at ville, at, naar et Menneske eller en Engel vil, at dette
eller hiint Legeme sættes i Bevægelse, det da bevæges."
% p. 22
Gud er altsaa, ifølge det anførte Exempel, Bevægelsens sande Aarsag, Engelens
Villie kun Anledningsaarsag.

„For at gjøre dette endnu mere indlysende, lad os antage, at Gud vilde det
Modsatte, af hvad nogle Aander, f.~Ex. Dæmonerne, ville; saa kunde man dog i
dette Tilfælde ikke sige, at Gud meddeler dem sin Magt, thi de kunne jo Intet
udrette af, hvad de ønske. Imidlertid vilde disse Aanders Villier dog være
naturlige Aarsager til Virkninger, de frembragte. Slige Legemer vilde da kun
bevæges til Høire, fordi disse Aander ønskede, at de skulde bevæges til Venstre:
de foranledige da Gud til at gjøre det Modsatte af, hvad de ville, og deres
Villier blive saaledes ogsaa Anledningsaarsager.\footnote{S.~392. Il n'y a donc
qu'un seul vrai Dieu et qu'une seule cause qui soit véritablement cause: et l'on
ne doit pas s'imaginer que ce qui précéde un effet en soit la véritable cause.
Dieu ne peut même communiquer sa puissance aux créatures, si nous suivons les
lumiéres de la raison: il n'en peut faire de véritables causes: il n'en peut
faire des dieux. Mais quand il le pouroit, nous ne pouvons concevoir pourquoi il
le voudroit. Corps, esprits, pures intelligences, tout cela ne peut rien. C'est
celui, qui a fait les esprits, qui les éclaire et qui les agite... Enfin c'est
l'Auteur de nôtre être qui exécute nos volontez, semel jussit, semel paret. Il
remue même nôtre bras lorsque nous nous en servons contre ses ordres, car il se
plaint par ses Prophetes que nous le faisons servir à nos desirs injustes et
criminels.}

\begin{center}---\end{center}

Der er i Theologiens Historie ingen Dogmatiker, som med større Strenghed har
fastholdt Begrebet om Guds Almagt, end Malebranche, Ingen, som med større
videnskabelig Conseqvens har udtalt, hvad der ligger i Skabelsesdogmet, end netop
Malebranche.

Han har ikke, som en Calvin, indskrænket sig til en logisk Form for et religiøst
Magtsprog; han har som philosophisk Tænker overalt stræbt at gjøre
Troessandheden tilgjængelig for Fornuften, og desaarsag heller ikke undslaaet sig
for at binde an med de virkelige Phænomener, og give os en med hans Theori
overeensstemmende Udtydning af den sandselig nærværende Tingenes Orden.

Vel har Malebranche ikke paa speculativ Maade kunnet see ind
% p. 23
i Faderens Urgrund, opdage Verdensspirer i Sønnens Pleroma, og mediere
Skabelsens Proces i Aandens plastiske Kraft; men Dogmet har han, det rene Dogme
om Skabelse af Intet. Han har om dette Dogme ikke blot erklæret, at det maa
tænkes (thi fordi det engang imellem siges om et Dogme, at det maa tænkes,
derfor er det jo ikke sagt, at det virkelig bliver tænkt); nei, han har i Alvor
forsøgt at tænke det, holde Tanken fast, og anvende den til en videnskabelig
Forstaaelse af alt det Skabte.

At anvende Skabelsestanken til en videnskabelig Forstaaelse af det Skabte, vil
sige, at see den hele Verden og hver enkelt Ting i Verden under det dobbelte
Synspunkt: \emph{Intet og dog virkelig Noget!}

Det Begreb, hvori denne Modsigelse opløses, er Begrebet: Almagt; det Skabte er i
sig et Intet, men paa Almagtens Bud er det blevet til Noget.

Den, der nu seer paa det Tilblevne saa eensidig, at han kun seer dets Intethed,
for ham er Verden Skin og Skygge. Om Theologien vil kalde ham Pantheist eller
Akosmist, er ligegyldigt; Sagen er, at han fornegter Skabelsen, idet han
fornegter Almagten.

Den, der seer paa det Tilblevne saa eensidig, at han kun seer dets Realitet,
glemmer dets Intethed, og tillægger det Væsenhed, for ham er Tilblivelsen en
Naturproces, og Verden et Naturens Rige. Om Dogmatiken forøvrigt maa henregne ham
blandt Deister, Naturalister eller Atheister, er i denne Sammenhæng ligegyldigt;
Sagen er, at han fornegter Skabelsen, idet han fornegter Almagten.

Den derimod, som i alle Forhold under alle Belysninger bestandig seer det
Tilblevne fra den dobbelte Side: som Intet i sig og dog som virkelig Noget, han
seer i Alt Guds Almagt, han tænker, om end kun menneskelig, Skabelsens Tanke, han
fastholder Skabelsens Dogme.

Det er dette, Malebranche har gjort; og man skal i Theologiens Bibliotheker
vistnok lede længe, inden man, hvad dette Punkt angaaer, finder en Tænker saa
ærlig conseqvent som ham.

Skabte, af Intet ere ham alle tilværende Ting i sig Intet, og vedblive
% p. 24
at være det; skabte af Almagten, ere ham de tilværende Ting virkelig til i en
virkelig Verden, som Skaberen opholder og styrer.

Hvad er Materien? Det Udstrakte, der synes at være Noget, naar det sees
sandselig; men, seet med Fornuftens Øie, er Materien væsenløs, uden Form, uden
Evne, uden Kraft til at virke paa nogensomhelst Maade i nogensomhelst Retning.
Desuagtet kan man, ifølge Theorien, ikke sige, at Materien ikke er til;
tvertimod! den er virkelig til; thi den er jo skabt: den almægtige Villies Alvor
indestaaer for dens Tilværen.

Hvad er Sjælen? Det Tænkende, der i endelige Tanker tænker paa sig selv, og
bilder sig ind at være Noget. Men Sjælen, den endelige Aand, er i sig selv et
Intet, et reent Intet, ligesaa uvirksom, død og magtesløs som Materien. Ikke
desto mindre er Sjælen virkelig til; den er jo skabt, og altsaa paa Almagtens
Bud af Intet bleven til Noget.

Som de skabte Tings Tilbliven er en mirakuløs Overgang fra Intet til Noget, er
ogsaa deres Tilværen en mirakuløs Svæven mellem Noget og Intet. Corps, esprits,
pures intelligences tout cela ne peut rien. C'est celui, qui a fait les esprits,
qui les éclaire et qui les agite.

Hvad der gjælder om Tingene selv, gjælder da tillige om deres indbyrdes Forhold.
Dette Forhold kan umulig være væsentligt, da Tingene selv ere uden Væsenhed. Her
er ingen Vexelvirkning imellem de tilværende Ting: de naturlige Aarsager ere jo
ikke sande Aarsager. Ikke desto mindre er her dog i denne virkelige Verden en
virkelig Sammenhæng, virkelige Love, en virkelig Tingenes Orden; ere de skabte
Ting end ikke sande Aarsager, saa ere de dog for den skabende Villie tilværende
Anledningsaarsager: de guddommelige Love ere Aabenbarelser af den høieste
Fornuft.\footnote{S.~390. Dieu a créé le monde... et produit ainsi tous les
effets que nous voyons arriver, parce qu'il a voulu aussi certaines loix, selon
lesquelles les mouvemens se communiquent à la rencontre des corps: et parce que
ces loix sont efficaces, elles agissent, et les corps ne peuvent agir.}
% p. 25
Man har nu rigtignok meent, at kunne gaae et Skridt videre end Malebranche, og
forbedre hans Skabelseslære ved at tillægge det Skabte dog nogen Selvstændighed,
noget Anlæg til indre Selvudvikling, nogen Evne til Virksomhed; det er da navnlig
Leibnitz, som her har villet corrigere eller rettere tilbagetrænge
„Mirakeltheorien" ved at opfinde Monadelæren, udtænke harmonia præstabilita, og
hævde principium rationis sufficientis.

Det er imidlertid, hvad Dogmet angaaer, en mislig Forbedring. Det Skabte skal da
ved Skabelsen ikke blot have faaet virkelig Tilværen, men med det Samme en
Tilværen i relativ Selvstændighed.

Det er nu denne relative Selvstændighed, Malebranche ikke har kunnet tænke sig,
og en Leibnitz, naar vi see nærmere til, nok ikke heller. Thi hvad er relativ
Selvstændighed? En Grad af Tilværen udenfor Almagten, en Selvvirken, som, om den
end for Resten er nok saa ringe, dog alligevel ikke er Almagtens Virken, altsaa:
en virkelig Grændse for Almagten.

Er det nu Alvor med denne virkelige Grændse, saa er jo Almagten en virkelig
begrændset Magt, og ikke længere Almagt. Har den Almægtige derimod endnu, som
før, Grændsen i sin Magt, saa er den virkelige Grændse jo endnu, som før, et
Intet i sig, og det mirakuløse Forhold mellem Skaber og Skabning maa vedblive.

Et Forhold, der er begyndt med et Mirakel, og grundlagt ved et Mirakel, maa ogsaa
fortsættes ved et uophørligt Mirakel: ikke blot i nogle Øieblikke, ikke blot, om
jeg saa maa sige, en 8 à 14 Dage efter Skabelsen; men i Aarhundreder, i
Aartusinder, i Evigheders Evighed maae de skabte Ting vedblive at være i sig,
hvad de oprindelig ere og have været, et Intet, et reent Intet for den Almægtige.

Det er denne Conseqvens, Malebranche fastholder, og derfor kan han ikke tænke sig
Selvvirksomhed i nogensomhelst Skabning, være sig Engel eller Menneske, thi
væsentlig Selvvirksomhed er en guddommelig Sag, og Gud selv kan ikke skabe Guder:
il n'en peut faire de véritables causes, il n'en peut faire des Dieux.
% p. 26
Den Forbedring, Leibnitz indfører, bestaaer nu deri, at han lader Gud skabe
Guder\footnote{Chaque esprit étant comme une petite divinité dans son
département. Monadol. --- Les esprits étants comme de petits Dieux. Syst.
nouv.}; thi hans relative Monader ere jo som en Vrimmel af Smaaguder. Enhver af
disse Smaaguder bærer paa Guders Viis en Verdens Mulighed i sit Indre, er paa
Guders Viis sig selv nok, fører paa Guders Viis et evigt Liv, og kan, om den end
svækkes til en given Tid, dog i al Evighed ikke (det skulde da være ved et
extraordinært Mirakel) tabe Selvudviklingens Evne\footnote{Atque ideo etiam nulla
datur generatio, nec mors perfecta rigorose loquendo. Sunt emim evolutiones et
accretiones, quas generationes appellamus, quemadmodum involutiones et
diminutiones, quod mortem vocamus. Princip. Philos. Thesis 76.}.

Men hvad er nu vundet ved at lægge al denne Selvvirksomhed ind i det Skabte? To
halve Tanker istedetfor een heel: Conseqvensen er brudt. Paa den ene Side skal
det endnu være Alvor med Skabelsen, altsaa med den Tanke, at Almagtens Grændse er
i sig et Intet. Paa den anden Side skal det tillige være Alvor med det Skabtes
Selvstændighed, og med den Tanke, at Almagtens Grændse er i sig et Noget.
Almagten skal altsaa paa een Gang være baade almægtig og dog ikke almægtig; see
det er Det, der er vundet.

Er nu Dette, hvad Leibnitz vil indføre, selv naar det gives en mere fattelig og
vulgær Form, virkelig et Fremskridt; ja, saa lader det sig rigtignok ikke negte,
at Theologien nødvendigviis maa have Noget at sige i det Videnskabelige; thi nu
bliver Forholdet mellem Gud og Verden saa forviklet, som muligt.

Snart har Almagten Grændse, og der skeer intet Mirakel; snart har Almagten ingen
Grændse, thi der er dog skeet et Mirakel. Her er Meget at overveie, naar det, vel
at mærke, skal overveies videnskabelig. Her er nihil negativum, og her er nihil
privativum; her er omnipotentia, her er præscientia, her er providentia, her er
conservatio; her er conservatio rerum simplicium, her er annihilatio rerum
simplicium, her er conservatio nexus
% p. 27
cosmici, her er destructio nexus cosmici; og saa er her desforuden concursus, og
concursus igjen baade universalis og specialis og specialissimus.

Kommer saa hertil en Mængde Skriftsteder, vel exegetiserede, vel citerede, vel
grupperede Skriftsteder; ere saa disse Skriftsteder i et passende Schema
ledsagede af kritisk-dogmatiske Bemærkninger; give saa disse Bemærkninger tillige
Vink om, at det med hine supranaturalistiske Former nok ikke har stort at betyde;
ere saa disse Vink i den Grad tydelige, at Enhver maa kunne see, hvor Grændsen
er, og hvor Grændsen ikke er, og følgelig selv indsee, baade hvad der maa tænkes,
og hvad der skal tænkes, og at her egentlig slet ikke tænkes: saa have vi et til
Sagens Vigtighed nogenlunde passende Indtryk af den Behandling, som det ved
„Skabningens Selvvirksomhed" indsmuglede Naturbegreb undergaaer i Theologien.

Ved at ville forbedre „Mirakeltheorien" taber Leibnitz sin videnskabelige
Holdning, og stiller sig i deres Række, der ellers qvakle med concursus og nexus
cosmicus og destructio nexus cosmici\footnote{Omnis autem monas est
inexstingvibilis, neque enim substantiæ simplices nisi \emph{creando vel
annihilando}, id est miraculose oriri aut desinere possunt. Epist. ad
Bierlingium.} o. desl.; kun at han gjør det med Talent og diplomatisk Smidighed.

Thi hvad vil det dog egentlig sige, at der til Hverdagsbrug ikke skee Mirakler,
men kun ved overordentlige Leiligheder? Det vil, ligefrem forstaaet, sige med
rene Ord: til Hverdagsbrug er Gud ikke almægtig; men ved overordentlige
Leiligheder er han almægtig.

At de skabte Ting ere i sig Intet, og deres naturlige Sammenhæng i sig Intet,
denne Tanke kunne vi Mennesker ikke udholde til Hverdagsbrug, den træder kun i
Kraft ved enkelte overordentlige Leiligheder; for nu ikke at overvældes af
Forestillingen om Guds Almagt, Alvidenhed og Allesteds%
% p. 28
nærværelse, maa man tilbagetrænge den altfor paatrængende Mirakeltheori, og bøde
paa Tingenes Realitet med lidt commercium substantiarum, harmonia præstabilita,
lidt nexus cosmicus, principium rationis og andre deslige Middelbarheder.

„Opholdelsen er en fortsat Skabelse"; dersom denne Tanke virkelig fastholdes, saa
har Malebranche Ret. Der kan da umulig være anden Forskjel mellem Skabelse og
Opholdelse end den, der ligger i de to Begreber: Begyndelse og Fortsættelse.

Den samme Almægtige, der fordum sagde: Bliv! maa endnu i hvert Secund, i hvert
Øieblik, blive ved med at sige: Vedbliv!

Dersom man nu desuagtet vil have en anden og „grundigere" Forskjel mellem
Skabelse og Opholdelse, hvor komme vi saa hen? Saa komme vi til det Resultat, at
Skabelsen er det umiddelbare, Opholdelsen det middelbare Mirakel. Men i denne
„Grundighed" er der ingen Mening; thi et middelbart Mirakel er en Modsigelse.
Dieu n'a pas besoin d'instrumens pour agir, il suffit qu'il veüille.

Theologien er for tilbøielig til at forblande Kategorierne: Eensidighed og
Conseqvens. Naar Tænkeren engang imellem griber en af dens egne Forudsætninger,
og udvikler, hvad deri væsentlig indeholdes, saa hedder det strax: hvilken
Eensidighed!

Slige Afviisninger ere uvidenskabelige. I Videnskaben har een heel Tanke, som
virkelig bliver tænkt, mere Værd end ti halve Tanker, som paa Grund af deres
indre Uforenelighed aldrig blive tænkte, hvormegen Umage man end gjør sig med at
forklare, at de „maae tænkes".

Malebranche er med Hensyn paa Skabelsesdogmet ikke en eensidig, men en conseqvent
Tænker. Det er ham Alvor med Skabelsen; det er ham derfor ogsaa Alvor med
Opholdelsen: han tænker ikke pro forma! Saalidt som det Tilblevne fordum har
kunnet skabe sig selv, saa lidt kan det endnu den Dag idag opholde sig selv; saa
lidt som det Tilblevne oprindelig har kunnet
% p. 29
skabes af et andet Tilblevet, saa lidt kan endnu det ene Tilblevne enten opholdes
eller bevæges af det andet\footnote{S.~391. Je dis de plus qu'il n'est pas
concevable, que Dieu puisse communiquer aux hommes ou aux Anges la puissance
qu'il a de remüer les corps; et que ceux qui prétendent, que le pouvoir que nous
avons de remüer nos bras, est une véritable puissance, doivent avoüer que Dieu
peut aussi donner aux esprits la puissance de créer, d'anéantir, de faire toutes
les choses possibles, en un mot, qu'il peut les rendre tout-puissans.}.

Ifølge Malebranche kan Skabelsesdogmet umulig beroe paa den Modsigelse, at Gud i
Tidernes Begyndelse skulde have anvendt sin Almagt til at sætte sig selv en
væsentlig Skranke, en Skranke, som er i sig et Noget, en Skranke, han ikke med
hvert Øieblik havde absolut og ubetinget i sin Magt; nei, Skabelsen kjendes
derpaa, at Guds Almagt af Kjærlighed og med Frihed sætter sig en Modsætning, og
igjennem denne Modsætning paa alle Punkter, i alle Retninger, til alle Tider
forholder sig til sig selv, sig selv alene.

Derfor er her en virkelig Verden, og denne Verdens Tilværen er Guds Værk; derfor
er her i Verden en virkelig Orden og virksomme Love, thi denne Orden er Guds
Viisdom, disse Love ere Guds Tanker. Derfor er det Alvor med Loven og med de
tilsvarende Anledningsaarsager, i den Grad Alvor, at Malebranche just seer Guds
Frihed deri, at den Almægtige, som i hvert Øieblik virker enhver Tilblivelse og
enhver Bevægelse, binder sig ved sine egne Tanker, sin egen Villie, sine egne
Love. Semel jussit, semel paret. Il remue même notre bras, lorsque nous nous en
servons contre ses ordres, car il se plaint par ses Prophetes que nous le faisons
servir à nos desirs injustes et criminels.

Sandheden i det Naturlige er ikke Natur; det Naturlige er ikke naturligt: det er
en Guds Gjerning.

Den gjennemførte Occasionaltheori kan betragtes som en videnskabelig Udlægning af
Psalmistens Ord (Ps. 104, 24): \emph{Herre, hvormange ere dine Gjerninger, Du
gjorde dem alle viselig.}
% p. 30
Vi have altsaa ikke gjort Theologien Uret ved at fremstille Malebranche som et
Exempel paa theologisk Videnskabelighed.

\begin{center}---\end{center}

Den i Theologien gjennemførte \emph{Skabelseslære} er en gjennemført
\emph{Naturfornegtelse}.

Det theologiske Naturbegreb er negativt; har dette negative Begreb videnskabelig
Gyldighed, saa er Naturlæren ikke blot en meget underordnet, men tillige meget
farlig Videnskab, thi Naturlærens Naturbegreb er positivt.

Hvad der er fatteligt, hvad der er umiddelbart indlysende, hvad der er
umiskjendeligt, hvad der er uimodsigeligt, om alt Saadant sige vi: „det er
naturligt". At en Steen, der kastes i Luften, maa falde til Jorden, er naturligt;
at Den, der holder sin Finger i Flammen, maa faae den brændt, er naturligt; at
Den, der falder i Vandet, maa blive vaad, Den, der ikke kan trække Veiret, qvæles
o. s. v., er Altsammen naturligt. Det maa være saaledes; det kan umulig være
anderledes. Hvorfor ikke? Man give saamange Forklaringer og føre saamange
Beviser, man vil: Nerven i alle Beviser, det Forklarende i alle Forklaringer, er
dog alligevel det Ene: „ja det er naturligt!"

Men i alt det Naturlige kan jo dog Intet være mere naturligt end Naturen; man maa
derfor ikke undres over, at Menneskeheden, trods det „høiere Lys", hvormed den
theologiske Videnskabelighed alt i Aarhundreder har lyst for Slægterne, dog endnu
den Dag idag saa uvilkaarlig fanges i den Indbildning, at Naturen virkelig er, at
der i Tilværelsen er Noget, vi i naturlig Forstand maa kalde Naturens Væsen.

Naturens Væsen! det vilde maaskee falde nok saa naturligt at sige: Natur og
Væsen; thi disse to Begreber ere saa inderlig beslægtede, at det ene i Grunden
maa forklare det andet. Enhver Begrundelse gaaer i Videnskaben ud paa at
eftervise, hvad der nødvendigviis maa følge af „Sagens Væsen", men hvad der følger
af „Sagens Væsen" er jo kun Det, der ligger i „Sagens Natur".
% p. 31
Ere nu begge Begreber i Grunden Eet, og er Grund og Begrundelse afgjørende for
al Videnskab; saa kan man da heller ikke undre sig over, at Menneskeaanden endnu
bestandig fastholder den Tanke, at der dog alligevel er og maa være Sandhed i en
Videnskab, som overalt gaaer ud paa at erkjende Naturen i Naturens Lys, uagtet
Naturlærens naturlige Grunde vistnok ere komne i en høist betænkelig Conflict med
Theologiens „dybere Grunde".

Skal Naturlæren gjælde for sand Videnskab; da maa der begyndes med den
Forudsætning, \emph{at der er et naturligt Vexelforhold imellem Sandserne og den
sandselige Verden.}

Men, dersom denne Forudsætning skal gjælde, maa Theologien see at finde sig en
Plads udenfor Videnskaben; thi dens Videnskabelighed, som Theologi, beroer jo
netop paa en Fornegtelse af denne Forudsætning.

Den, der vil erkjende Naturen, maa stole paa Sandsernes Vidnesbyrd. At stole paa
Sandsernes Vidnesbyrd, er ikke det Samme, som at bedømme Gjenstandene efter det
første det bedste umiddelbare Sandseindtryk. Hvorvidt Naturlæren i denne
Henseende forstaaer at gjøre Forskjel mellem Skin og Væsen, har den dog vel
blandt Andet tilstrækkelig viist i sit Forsvar for det copernikanske System.

At stole paa Sandsernes Vidnesbyrd, vil sige, at stole paa den indre væsentlige
Overeensstemmelse mellem det Sandselige og Det, der sandses, det er: paa
Sandseindtrykkets naturlige Sandhed.

„Det er en Sandhed, at Sandserne sandse", dermed maa der begyndes i Naturlæren;
„det er en Usandhed, at Sandserne sandse", dermed maa der begyndes i Theologien.

Det er, videnskabelig forstaaet, ikke en Eensidighed mod en Eensidighed, vi her
have for os; det er Princip mod Princip, Videnskab mod Videnskab, hvorom det
gjælder.

Ved at stole paa Sandserne kan et Menneske, selv med den bedste Villie, aldrig
bringe det videre end til at see Naturen i et naturligt Lys.
% p. 32
Jo mere Sandserne gjælde, desto mindre gjælder Troen; jo mindre Sandserne
gjælde, desto mere gjælder Troen.

Der er en uendelig Forskjel imellem den Tillid til Sandsernes og Erfaringernes
Vidnesbyrd, der kræves i Naturlæren, og den Tillid til den Helligaands
videnskabelige Bistand, Theologerne kalde Tro.

Vil man desuagtet udvide Begrebet: Tro, til at omfatte enhver Antagelse, der
beroer paa inderlig Stolen og Tillid, saa kan man dog vel ikke negte, at her er
Forskjel mellem Tro og Tro. Naturforskerens Tro er Naturbekræftelse a priori;
Theologens Tro er Naturfornegtelse a priori.

I Videnskaben (her glemmes ikke, at det er i Videnskaben) maa enten den naturlige
Tro tilintetgjøre den overnaturlige, eller omvendt.

At ville see Naturen paa een Gang \emph{baade} i et naturligt og dog
overnaturligt Lys, \emph{baade} i et helligt og dog videnskabeligt Lys, denne
Villen-baade-og er ikke Tegn paa speculativ Dybsindighed, men et Tegn paa
Tankeløshed; thi det er en Modsigelse.

Skal det naturlige Lys skinne, maa det overnaturlige slukkes; skal det hellige
Lys oprinde, maa det videnskabelige slukkes.

C'est pour nous accommoder à la maniére ordinaire de parler, que nous dirons dans
la suite, que les sens sentent; med disse Ord slukker Malebranche det naturlige
Lys; den fromme Tænker „skuer alle Ting i Gud": fides præcedit intellectum.

Har Theologien nu virkelig en saadan Tillid til sit: fides præcedit intellectum,
at den dermed i vore Dage for Alvor tør byde Naturlæren Spidsen; ja da skulle vi
ikke blot tilstaae, at den har en Tro, men endog villig indrømme, at dens Tro
gaaer langt forud for dens Forstand.

Skal Naturlæren ansees for sand Videnskab, da maa den Grundsætning fastholdes,
\emph{at den naturlige Forstand i Sandhed formaaer at erkjende Naturens Love.}

Men, dersom denne Grundsætning skal gjælde i Videnskaben, gjør
% p. 33
Theologien vistnok bedst i at søge sig en Plads udenfor Videnskaben; thi det er
en Grundsætning, Troen maa fornegte.

Skal Troen dømme, saa ere Tilværelsens Love Guds Tanker, hans Villie den eneste
sande Aarsag. Il n'y a donc qu'un seul vrai Dieu et qu'une seule cause, qui soit
véritablement cause.

Skal Forstanden dømme, da ere Tilværelsens Love grundede i Tingenes Væsen, og der
maa følgelig være ligesaa mange sande Aarsager, som der ere naturlige Kræfter.

Skal Troen dømme, maae Tingenes Love kunne forandres, hvad Øieblik det behager
Gud; skal Forstanden dømme, maa endog den mindste Afvigelse fra Loven antages at
stride imod „Sagens Natur".

Her gaaer Dom imod Dom, og det gjælder Naturens Væsen. Naar den ene Dom høres,
kan den anden ikke høres; man kan ikke paa een Gang høre to Domme, fordi man har
to Øren.

Troens Øre kan slet ikke høre Forstandens Dom, saalidt som Forstandens Øre kan
høre Troens Dom. Troens Dom høres i „Oratoriet", men Forstandens Dom maa høres i
„Laboratoriet". At stille sig midt imellem „Oratoriet" og „Laboratoriet", spidse
begge Øren, og sige: credo, ut intelligam, for at lytte til begge Sider, er at
give sig selv une fausse position.

Skal Forstandens Dom høres i Videnskaben, maa Troens Dom høres udenfor
Videnskaben; man kan ikke forlige saa contradictoriske Domme i een videnskabelig
Dom.

For at dølge den skjærende Modsigelse, har man forsøgt en human Accommodation.

„Naturens Love ere visselig Guds Tanker; men Gud kan jo ikke forandre disse
Tanker, da det strider imod hans Uforanderlighed".

Saa lyder den humane Accommodation; men Sandheden trænger paa: det Ueensartede
vil adskilles.

Det gaaer med human Accommodation til en Tid, som med jesuitisk
% p. 34
Tortur til en Tid; pludselig kommer --- Øieblikket! Aanderne vredes, Principerne
harmes: Luther er grov, Galilæi stamper i Gulvet.

Skal Naturlæren ansees for en Videnskab, da gjælder den Fornuftidee, \emph{at
Naturriget er et i sig sammensluttet Hele, hvis ueensartede Grunddele ifølge
deres Qvaliteter forbindes og sondres saaledes, at Heelhedens særegne Led formes,
dannes og opløses med indre Nødvendighed.}

Men, dersom denne Fornuftidee skal gjælde i Videnskaben, gjør Theologien bedst i
at flytte ud af Videnskaben; thi det ligger jo netop i dens supranaturalistiske
Charakteer, at fornegte denne Idee.

„Grund", „Urgrund", „Vorden", „Fødsel", „indre Formen", „uvilkaarlig Yttring af
eget Anlæg,"\footnote{„Unsere Erklärung von dem Natürlichen", siger Link
(Propyläen der Naturgeschichte. Berlin 1839, S.~4) ist dieselbe, welche
Aristotoles giebt," og citerer i denne Anledning Physic. lib. 2, c. 1.

Τὰ μὲν γὰρ φύσει ὄντα πάντα φαίνεται ἔχοντα ἐν ἑαυτοῖς ἀρχὴν κινήσεος καὶ στάσεος,
τὰ μὲν κατὰ τόπον, τὰ δὲ κατ' αὔξησιν καὶ φθίσιν, τὰ δὲ κατ' ἀλλοίωσιν.} „Instinct"
og „Drift", ere de Kategorier, der foresvæve Bevidstheden, naar Talen er om den
moderlige Natur.

At der er en Sandhed heri, anerkjendte Theologien da ogsaa paa sin Viis, dengang
den ret begyndte at speculere over den faderlige Natur.

I Læren om Sønnens evige Fødsel af Faderen fik Athanasius jo Ret imod Arius, og
hvorfor? Fordi Fornuften i dette Punkt fik Lov til at indsee, at, dersom Sønnen
var i Væsenseenhed (ὁμοούσιος) med Faderen, saa kunde han ikke være skabt med
Frihed, men maatte være født af Nødvendighed.

Men, idet Theologien herved gjorde et Indblik i Naturens Væsen, bestod dens
videnskabelige Fremskridt egentlig deri, at den overførte Kategorier fra det
Naturlige paa det Overnaturlige. I sin Stræben efter at begribe den overnaturlige
Natur (θεία φύσις) spærrede den sig da selv Veien til en Forstaaelse af den
naturlige Natur.
% p. 35
Sønnens Natur er født, men ikke skabt; Adams Natur er skabt, men ikke født; Adams
Børn ere derimod baade skabte og fødte: kan der komme noget Videnskabeligt ud af
Dette?

Gud kan ikke skabe Guder, derfor er Sønnen født fra Evighed af; det er Det,
Athanasius har indseet.

Gud kan ikke skabe Guder, derfor er der ingen Natur i det Skabte; det er Det,
Malebranche har indseet.

Hvad der i Sandhed skal blive til en Natur, maa blive til paa en naturlig Maade;
hvad der er blevet til ved et Mirakel, maa have en mirakuløs Tilværen, og kan
aldrig blive til noget i Sandhed Naturligt, saa lidt som det nogensinde kan blive
naturlig sand Aarsag til andet Naturligt. „I de naturlige Naturforhold er der
ligesaa lidt Sandhed, som i den naturlige Natur"; dette er den
theologisk-videnskabelige Conseqvens af Dogmet om Skabelsen.

Er Adam skabt, saa er Adam ingen Natur, saa er der ingen Natur i Adam. Er der
ingen Natur i Adam; saa er der heller ingen Natur i Eva. Er der ingen Natur i
Adam og Eva, saa kunne Adams og Evas Børn heller ikke fødes; de maae alle skabes.
I denne Forstand kan man da med Sandhed sige, at hvert Menneske er en Guds
Skabning.

Menneskelig Natur, menneskelig Fødsel, alt Sligt er, ifølge Malebranche,
ligesaavist Skin og Sandsebedrag, som det er Skin og Sandsebedrag, at det ene
Legeme virkelig skulde være istand til at sætte det andet i Bevægelse.

At et Menneske kan fødes, lader sig tænke; det er en reen, naturlig Sag, og der
er forsaavidt Eenhed i Tanken. At et Menneske kan skabes, lader sig tænke; det er
en reen, overnaturlig Sag, og der er forsaavidt Eenhed i Tanken. Men at et
Menneske skal kunne \emph{baade} fødes og skabes, \emph{baade} skabes og fødes,
dette Baade-og lader sig ikke tænke, thi det er en Modsigelse.

Dersom „Creatianisme" og „Traducianisme" virkelig lade sig mediere ved credo, ut
intelligam, saa lade Cirkel og Samvittighed sig nok ogsaa
% p. 36
mediere ved credo, ut intelligam. En „Fødselsproces", som tillige er „en
Skabelsesproces", en Fødselsskabelse, en Skabelsesfødsel, en traduciansk
Creatianismus, en creatiansk Traducianismus, en naturlig Mirakelproces, en
mirakuløs Naturproces: dersom dette er videnskabeligt, saa er Naturlæren
uvidenskabelig.

Enhver, som veed, hvad Videnskab er, maa dog virkelig tilstaae, at Nutidens
Theologi befinder sig i en egen Stilling. I samme Grad den nærmer sig det
Overnaturlige, maa den fjerne sig fra Naturen; i samme Grad den vil nærme sig en
Naturerkjendelse, maa den fornegte det Overnaturlige.

Paamindede, især af Kant, gik Theologerne en Tidlang frem efter den Grundsætning,
at det for Alting gjaldt om at holde sig til Naturloven og den praktiske Fornuft.

„Kunde vi blot paa en videnskabelig Maade faae de mange ufornuftige Mirakler
bort!" tænkte Rationalisterne, og saa begyndte de at bortforklare de mange
„ufornuftige" Mirakler. Det lykkedes, endog imod Forventning; det ene Mirakel
forsvandt efter det andet --- paa en videnskabelig Maade. Tilsidst blev kun eet
Mirakel tilbage, Skabelsen nemlig; dette Mirakel maatte man vel i det Enkelte
omforklare, men i det Hele ikke bortforklare; thi Skabelsens Mirakel var dog i
Grunden et ret fornuftigt Mirakel. Ved denne fornuftige Fremgangsmaade blev
Theologien unegtelig meer og meer fornuftig, meer og meer videnskabelig; tilsidst
manglede der kun Eet: Conseqvensen.

I Videnskaben kan en Conseqvens vel undertiden lade vente paa sig, men aldrig
udeblive. Den Strausiske Kritik var Rationalismens Conseqvens; Conseqvensen tog
den sidste Rest af det sidste Mirakel bort, og Theologien gik bagvendt ud af
Videnskaben.

Tugtet af Kritikens Svøbe, har Theologien igjen saa smaat begyndt med credo, ut
intelligam, at vende sig mod det Overnaturlige; vi kunne ikke Andet end agte paa
dens Omvendelse. Men en sand Omvendelse har ogsaa sin Conseqvens: „viis mig din
Tro af dine Gjerninger!"

Dersom Theologien virkelig gjør Alvor af at vende sig mod det Over%
% p. 37
naturlige, og gjennemfører Omvendelsens Conseqvens; da skal den aldrig mere grue
for Miraklet, aldrig mere bøie sig for en Naturlov.

Dersom Theologien virkelig gjør Alvor af at vende sig mod det Overnaturlige, og
gjennemfører Omvendelsens Conseqvens; da vil dens Kræfter voxe, da vil den blive
stærk: stærk i Credo, stærk i Intelligam, stærk i Miraklet! saa stærk i Miraklet,
at dens Mirakelstyrke slaaer „Alnaturen" ihjel, og gaaer seierrig --- ud af
Videnskaben.

Dersom Theologien virkelig gjør Alvor af at vende sig mod det Overnaturlige, og
gjennemfører Omvendelsens Conseqvens, da faaer den Fred for al Kritik; da samles
den til sine Fædre, og lander i en bedre Havn: Athanasius har længst beredt den
Sted; Malebranche staaer og venter i „Oratoriet".

\begin{center}---\end{center}

Imellem Theologien og Naturlæren er der ingen Forstaaelse: Theologiens
Naturbegreb er ikke naturligt, og Naturlærens ikke theologisk.

\begin{center}---\end{center}
% p. 38  (Indbydelsesskrift appendix: the graduates' vitae)
De trende Videnskabsmænd, som i det nu forløbne Aar have erhvervet academiske
Grader, have meddeelt følgende Optegnelser om deres Levnetsløb:

\begin{center}\textbf{1.}\end{center}

Jeg, \textbf{Martin Salomonsen}, er født den 9de Marts 1814 i Kjøbenhavn, hvor min
Fader, \emph{Samuel Salomonsen}, der nylig er død i en høi Alder, i en Række af
Aar har været Kjøbmand; min kjærlige Moder, \emph{Amalie Henriques}, har jeg alt
mistet for over 20 Aar siden. Min Skoledannelse erholdt jeg i Borgerdydskolen i
Kjøbenhavn, hvis Rector, Prof. M. Nielsen, med Omhyggelighed stræbte at uddanne
mig og afsendte mig i October 1832 med et anbefalende Vidnesbyrd til
Universitetet, hvor jeg bestod Examen artium med Laudabilis c. encom. publ.
Efterat jeg Aaret derpaa havde fuldendt den anden Examen med eenstemmig Laud. præ
cæt. begyndte jeg at lægge mig efter Lægevidenskaben, som jeg fra mine tidligere
Aar altid har følt Lyst og Kald til.

Skjøndt Naturvidenskaberne dengang næsten ikke foredroges for de medicinske
Studerende, indsaa jeg dog snart deres Betydning for vort Studium og anvendte
derfor i de første Aar en ikke ringe Tid paa disse Discipliner, der i H. C.
Ørsted, Schouw, Reinhardt og Forchhammer havde ligesaa udmærkede som berømte
Repræsentanter. Samtidigen hørte jeg flere Forelæsninger især over Anatomie og
Chirurgie paa det chirurgiske Academie hos Withusen, C. Fenger og Stein, medens
Herholdt, Saxtorph, Otto og navnligen Eschricht og O. Bang vare mine Lærere ved
Universitetet. Da jeg
% p. 39
imidlertid stedse har havt meest Interesse for den practiske Side af
Lægevidenskaben, besøgte jeg fra 1835 dagligen Frederiks Hospital, hvor jeg under
Veiledning af G. Møller og O. Bang, der altid har viist mig en faderlig Velvillie,
fik Leilighed til stadigen at iagttage en Mængde forskjellige Syge og derved paa
rette Maade blev indviet til min fremtidige Virksomhed.

Saaledes forberedt indstillede jeg mig i Foraaret 1838 til den forenede
medico-chirurgiske Examen, som jeg bestod med første Characteer. Imidlertid
vedblev jeg endnu i de første Aar derefter at besøge Forelæsningerne ved
Universitetet især over chirurgisk Anatomie, Physiologie og Accouchement, i
hvilket Fag jeg 1839 endvidere søgte at uddanne mig paa Fødselsstiftelsen under
S. Saxtorph.

Et Tilbud om som Læge at ledsage en Patient til Udlandet modtog jeg Aaret derpaa
i det Haab derved tillige at faae Leilighed til at udvide mine Kundskaber. Jeg
forlod Hjemmet i Høsten 1840, besøgte først Badene ved Ems og drog derpaa gjennem
Schweitz og Norditalien til Toscana, hvor jeg tilbragde en længere Tid, hørte en
Deel Forelæsninger i Pisa og Florentz og overværede paa begge Steder især de
cliniske Øvelser paa Hospitalerne; paa Tilbagereisen over Østerrig og Preussen
opholdt jeg mig nogen Tid i Wien.

Ved min Hjemkomst i Sommeren 1841 blev jeg ansat som Candidat paa Frederiks
Hospital, og her fungerede jeg i 3 Aar saavel paa den chirurgiske som paa den
medicinske Afdeling under G. Møller, O. Bang og S. Trier. Kort efter at min
Tjenestetid paa Hospitalet var udløben, blev jeg udnævnt til Districtslæge i
Kjøbenhavns 9de Hoveddistrict, et Embede jeg nu i 11 Aar har beklædt og da jeg
allerede dengang havde lagt Grunden til en practisk Virksomhed i min Fødeby,
ægtede jeg den 24de November 1844 en Slægtning af mig, \emph{Emma Henriques}, der
har beredet mig et lykkeligt Hjem og skjænket mig en elsket Søn.

Efterat jeg i Aaret 1848 var bleven Medlem af det kongelige medicinske Selskab og
1850 af Sundhedscollegiet kaldet til Censor i Physiologie ved
% p. 40
den medico-chirurgiske Examen for Bienniet 1850 og 51, erholdt jeg i Aaret 1854
af det medicinske Facultet Tilladelse til at disputere for Doctorgraden ved
Universitetet og den 9de December samme Aar forsvarede jeg paa Latin min paa
Modersmaalet skrevne Afhandling, „Udsigt over Kjøbenhavns Epidemier i den sidste
Halvdeel af det attende Aarhundrede".

Foruden dette Skrift har jeg næsten ikke udgivet Noget i Trykken; dog findes der
nogle casuistiske Meddelelser og en Comitebetænkning af mig i „Philiatriens"
Forhandlinger og som Medlem af det medicinske Selskabs hygieiniske Udvalg har jeg
fra Aaret 1850 først alene og siden i Forbindelse med Hr. Reservechirurg Paulli
givet en aarlig Udsigt over Kjøbenhavns epidemiske Constitution, der findes
aftrykt i Bibliothek for Læger.

\begin{center}---\end{center}

\begin{center}\textbf{2.}\end{center}

Jeg, \textbf{Carl Ludvig Müller}, er født i Kiøbenhavn den 9de Juni 1809. Min
Fader var \emph{Peter Erasmus Müller}, Professor i Theologien, min Moder
\emph{Louise Stub}, Datter af Commandeur-Capitain Stub. Jeg blev i mit 7de Aar
sat i Efterslægt-Selskabets Realskole og 3 Aar efter i Metropolitanskolen,
hvorfra jeg blev dimittert til Universitetet i 1826. Til Examen artium var jeg saa
heldig at blive den første blandt de Indkaldte og tog det følgende Aar ligeledes
anden Examen med Udmærkelse. Efter min Faders Ønske, ikke ifølge egen
Tilbøielighed, bestemte jeg mig for den theologiske Bane. Imidlertid følte jeg mig
snart tiltrukket ved de Studeringer af philologisk, historisk og philosophisk
Indhold, i hvilke Theologien førte mig ind, saavelsom ved de Forelæsninger jeg
hørte hos min Fader, Clausen og Hohlenberg. Jeg havde ikke nogen Grund til at
stræbe efter at tage Examen paa den kortest mulige Tid; derfor studerede jeg ikke
saameget
% p. 41
med dette Maal for Øie, men søgte, idet jeg ogsaa gjorde mig bekjendt med tydske
Theologers Skrifter, at komme til selvstændige Anskuelser, og anvendte længere
Tid paa de Discipliner, der meest interesserede mig. Da jeg i Sommeren 1832 tog
den theologiske Attestats erholdt jeg Charakteren Laud. et qu. egregie. De
paafølgende Aar forblev jeg i Hovedstaden, beskjæftiget med at fortsætte mine
Studeringer og tillige med at give Manuduction i Theologien. I 1833 foretog jeg
at besvare det ved Universitetet i den orientalske Philologi udsatte
Priisspørgsmaal, der gik ud paa en kritisk Udvikling af de Forandringer, det
hebraiske Sprog var undergaaet i det Tidsrum det Gamle Testaments Skrifter
omfatte, hvorved jeg fik Anledning til at udvide mine Kundskaber i de semitiske
Sprog. Prisen blev mig tilkjendt for denne Afhandling. Det følgende Aar led jeg
et smerteligt Tab, da Døden bortrev min høit agtede og elskede Fader, efterat han
kun i 4 Aar havde beklædt Sjellands Bispestol. I den Hensigt at udarbeide en
Afhandling for Licentiat-Graden i Theologien vendte jeg mig til Studiet af
Kirkefædrene og den gamle christelige Dogmehistorie. Den Disputats, som udgik
heraf, blev fuldendt i Begyndelsen af 1836 og havde til Titel: De resurrectione
Jesu Christi, vita eam excipiente et ascensu in coelum sententiæ quæ in ecclesia
christiana ad finem usque seculi sexti viguerunt. Den blev antaget af det
theologiske Facultet og forsvart den 15de Sept. f. Aar, hvorefter jeg ved
Reformations-Jubelfesten i Oktober erholdt den academiske Grad. Det hørte til et
af mine største Ønsker at komme til at foretage en længere Udenlandsreise;
Stipendier til en saadan vare blevne mig tilstaaede paa 2 Aar baade fra
Universitetet og fra Fonden ad usus publicos, og jeg begav mig paa Reisen kort
efter Disputationen.

Her indtraadte et Vendepunkt i mit Liv. Jeg havde fra Begyndelsen ikke havt nogen
Lyst til den geistlige Virksomhed, og en saadan var heller ikke bleven vakt hos
mig ved de theologiske Studeringer, der ikke fra den praktiske Side havde draget
mig til sig. Jeg havde derfor oftere været nærved at forlade den theologiske Vei;
dette blev paa Reisen til en fast Beslutning hos
% p. 42
mig. Under mit henved treaarige Ophold i Tydskland, Frankrig og Italien
beskjæftigede jeg mig fornemmeligt med Kunst og Antiquiteter, for hvilke jeg fra
min tidligere Ungdom havde havt særdeles Interesse; de store Museer i de
forskjellige Hovedstæder, især i Italien, frembøde et rigt Stof hertil; tillige
studerede jeg det arabiske Sprog, hvortil jeg især benyttede Opholdet i Jena og
Paris. Efter min Hjemkomst i 1839 bestemte jeg mig til at gjøre Kunst og
Archæologi til Gjenstande for mine Studier, for derved om muligt at bane mig en ny
Vei, men var uvis, i hvilken Retning jeg først skulde vende mig; da opmuntrede
Kong \emph{Christian} VIII og vor bekjendte Archæolog \emph{Brøndsted} mig til,
specielt at lægge mig efter Numismatiken. Jeg tog dette tilfølge, og gav mig
ifærd med at gjennemgaae Samlingen i det kongelige Myntcabinet, hvor ikke blot
Directeuren, Brøndsted, ved sine lærde Kundskaber, men ogsaa Etatsraad
\emph{Thomsen}, der var Inspecteur ved Cabinettet, ved sin grundige Indsigt og
praktiske Kyndighed var mig til særdeles Nytte. I 1841 blev jeg adjungert som
Inspecteur ved Myntcabinettet og virkelig Inspecteur i 1842, da Brøndsted var død
og Thomsen havde erholdt dennes Plads som Directeur. Da der i Slutningen af 1841
var indløbet den Efterretning fra Rom, at min yngre Broder, Maleren Adam Müller,
som opholdt sig der, led af en Brystsyge, der truede hans Liv, og min Moder og
Søster ønskede at reise til ham, ledsagede jeg dem paa denne Reise; min Broder
overstod dette Anfald af sin Sygdom, som først 3 Aar senere i Hjemmet endte hans
meget lovende Kunstnerbane; for mig medførte Reisen et 4 Maaneders Ophold i Rom,
hvilket jeg benyttede til at udvide mine Kundskaber paa Oldtidens og Kunstens
Gebet.

Under mine gjentagne Ophold i Rom havde jeg havt megen Leilighed til at være
sammen med \emph{Thorvaldsen}; i Kjøbenhavn var jeg ham behjelpelig med at ordne
de Antiquiteter, han havde omkring sig i sin Bolig, og da en stor Deel af hans
Værker og de ham tilhørende Kunstgjenstande, bestemte til det Museum, der var
under Bygning, var bleven midlertidigt opstilt paa Christiansborg Slot, paatog jeg
mig Opsynet over denne Samling, der blev
% p. 43
meget besøgt, især af Reisende. Efter Thorvaldsens Død i 1844, da alle de
Kunstværker og Samlinger, som han havde testamentert til Kjøbenhavn, vare blevne
Communens Eiendom, overdroges det mig at udarbeide en Beskrivelse af
Antiquiteterne. Disse bestode i betydelige Samlinger af forskjellig Slags, baade
af ægyptiske, græske og romerske Antiquiteter, foruden af Mynter, ogsaa af
Gemmer, Vaser o. s. v., hvilket bevirkede, at jeg maatte gjøre specielle Studier i
flere forskjellige Retninger, og at jeg foretog en ny Reise til Berlin, for at
gjennemgaae det derværende rige Museum efter dets videnskabeligt udarbeidede
Cataloger. Da Thorvaldsens Museum i 1847 var opført, bleve
Antiquitets-Samlingerne af mig opstilte og ordnede i de Værelser og Skabe, som
dertil vare bestemte, og det følgende Aar blev jeg valgt til Inspecteur for
Museet af den af Thorvaldsen udnævnte Direction i Forbindelse med en af
Communalbestyrelsen nedsat Comitee. Jeg ordnede derpaa ligeledes Samlingerne af
Tegninger, Kobbere, Bøger m. m., og forfattede, foruden Beskrivelsen over
Antiquiteterne, der udkom i 3 Bind, ogsaa en Hovedcatalog over Museets øvrige
Afdelinger, som udgjør 5 Hefter. Beskrivelsen over Samlingerne i Thorvaldsens
Museum blev udgivet baade paa dansk og fransk i Løbet af Aarene 1847--1851.

I det kongelige Myntcabinet havde jeg overtaget de Afdelinger, der indeholde de
antike og de orientalske Mynter, ved hvilke sidste det Kjendskab til det arabiske
Sprog, jeg tidligere havde erhvervet mig, kom mig til Nytte. I Sommeren 1852
gjorde jeg en Reise igjennem Holland og Belgien til London og Paris og besøgte i
det følgende Aar Tydsklands forskjellige Hovedstæder, især i det Øiemed at
undersøge de store Myntsamlinger, knytte Forbindelser med Directeurerne for disse
og andre Numismatikere, og samle Materialier til videnskabelige Arbeider. En Deel
af det indsamlede Stof behandlede jeg i en Monographi: „Den macedoniske Konge
Philip IIs Mynter", hvilken jeg indgav til Universitetet som Dissertation for
Doctorgraden i Philosophien; den blev antaget og udkom i Trykken i Begyndelsen af
dette Aar; ifølge Forslag af det philosophiske Facultet blev jeg ved en kongelig
Resolution
% p. 44
fritaget for det mundtlige Forsvar af samme. Samtidigt havde jeg udarbeidet et
Værk over Alexander den Stores Mynter; dette, til hvis Udgivelse det kgl.
Videnskabernes Selskab har bidraget, er færdigt fra Pressen og udkommer i disse
Dage paa fransk under Titelen: Numismatique d'Alexandre le Grand.

I Aaret 1839 ægtede jeg \emph{Vilhelmine Levetzau}, Datter af afdøde Kammerjunker
C. Levetzau; dette Ægteskab, som udgjør mit Livs Lykke, har skjænket mig 4 Sønner
og en Datter.

I 1847 blev jeg udnævnt til Ridder af Danebrogen; i 1853 erholdt jeg Titel af
Professor.

\begin{center}---\end{center}

\begin{center}\textbf{3.}\end{center}

Jeg, \textbf{Nicolaus Ludvig Helweg}, er født den 26de April 1818 i Odense, hvor
mine Forældre Distriktslæge, Dr. \emph{Hans Zacharias Helweg} og \emph{Anna
Elsabeth født Wendt}, Aaret iforveien havde taget Ophold. Efter at have
gjennemgaaet den derværende Skole blev jeg i Efteraaret 1835 dimitteret til
Universitetet, og underkastede mig i Mai 1840 den theologiske Embedsexamen med
laud. I mine Studenteraar havde jeg fortrinsviis lagt mig efter det Gamle
Testamentes Studium, og indgav i Efteraaret 1839 en Besvarelse af den akademiske
Priisopgave i den orientalske Philologie, hvilken Afhandling tilkjendtes
„accessit". Kort efter forlod jeg Kjøbenhavn, og tilbragte tre Aar som Huuslærer
paa en Herregaard i Fyen; her sysselsatte jeg mig meget med den nyere tydske
Philosophie, især Schleiermacher og Hegel, og benyttede Udbyttet af disse Studier
i min første theologiske Afhandling „Om Tro og Viden", der blev trykt 1843 i 1ste
Bind af Tidsskrift „For Litteratur og Kritik, udgivet af Fyens Stifts litterære
Selskab". Samme Efteraar blev jeg af det nævnte Selskab valgt til Tidsskriftets
Redaktør, og besørgede denne
% p. 45
Forretning i fem Aar fra Kjøbenhavn af, hvor jeg i Foraaret 1844 atter tog mit
Ophold. Her sluttede jeg mig snart nøie til Grundtvig, der ogsaa tidligere som
Præst havde tiltalt mig meest, og blev under Studiet af Kirkefædrene og i det Hele
af Kirkens Historie altid mere vundet for hans eiendommelige, dristige og rige
Livsanskuelse. Jeg søgte at udvikle og begrunde denne min Betragtning fornemmelig
i tvende Afhandlinger: „Om Grundtvig og Schleiermacher" („For Litteratur og Kritik
6te Bind) og „Om Troen og Guds Ord, en exegetisk Undersøgelse" (Fortsættelser fra
Pedersborg udgivet af Dr. P. C. Kierkegaard 1ste Bind). Denne sidste havde været
indsendt til det theologiske Facultet som Disputats for Licentiat-Graden, men var
ikke bleven antaget; „Facultetet miskjendte ikke Arbeidets gode Sider", men fandt
de deri udtalte Paastande stridende „imod vor Kirkes Princip og imod den nyere
Videnskab". Til denne Tid udgav jeg hefteviis i Løbet af et Par Aar „Den Danske
Psalmedigtning" i to Bind (1846--47) i Forening med Cand. C. J. Brandt, med hvem
jeg ogsaa nu udgiver „Dansk Kirketidende". Tvende Gange har jeg søgt en akademisk
Lærerpost, hvortil jeg altid har følt Lyst og Kald; den ene Gang var ved
Universitetet i Christiania, og den norske Regjering kom mig strax med stor
Forekommenhed imøde, men senere forandrede Forholdene sig, og jeg fandt mig
foranlediget til at træde tilbage. Da der indtraadte en Vacance i det herværende
theologiske Facultet, søgte jeg denne Plads, og kom derved til i Foraaret 1852 at
holde fem theologiske Concurrence-Forelæsninger, men min Medconcurrent, nuværende
Prof. Lic. Hagen blev dengang foretrukken. --- I September 1847 indtraadte jeg i
Ægteskab med \emph{Caroline Thalia Elisabeth Luun}, Datter af Proprietær A. Luun
til Lammehauge i Fyen. Uden nogen offentlig Ansættelse har jeg i en Række af Aar
sysselsat mig med litterære Arbeider, og siden jeg i Efteraaret 1852 erholdt en
Portion af det Smithske Stipendium ganske anvendt min Tid til theologiske og
historiske Studier. „Den Danske Kirkes Historie efter Reformationen", der er
udkommet i Hefter siden 1849, vil snart fuldstændig foreligge i to Bind. Min
Afhandling om „De Danske Domkapitler" --- et For%
% p. 46
søg til en ny Behandling af den ældre Kirkehistorie --- blev af det philosophiske
Facultet kjendt værdig til at tjene som Disputats for Doctorgraden. Ved den
mundtlige Disputats den 11te Sept. d. A. blev efter min derom indgivne Begjæring
for første Gang ved vort Universitet Modersmaalet benyttet ved en saadan akademisk
Forhandling.

\begin{center}---\end{center}

Ved nærværende Skrift indledes Universitetets forestaaende Aarsfest til Erindring
om Kirkens Reformation, hvilken Fest er berammet til Torsdagen den 8de November
Kl. 12, i Universitetssalen.

Talen holdes af Prof. R. Nielsen.

Derpaa vil Universitetets Rector, \emph{J. E. Larsen}, meddele en kort Beretning
om de academiske Begivenheder i det nu forløbne Rectoratsaar, samt proclamere som
Doctor i det lægevidenskabelige Facultet: Cand. med. \emph{M. Salomonsen}, og som
Doctorer i det philosophiske Facultet: Professor, Licent. theol. \emph{C. L.
Müller} og Cand. theol. \emph{N. L. Helweg}.

Med denne Fest fratræder Universitetets nuværende Rector ligesom ogsaa
Faculteternes nuværende Decaner. Rectoratet overtages af \emph{J. N. Madvig} som
valgt Rector for det forestaaende Aar, og Decanatet i de forskjellige Faculteter
gaaer over til de nye valgte Decaner, nemlig i det theologiske Facultet:
\emph{C. E. Scharling}; i det rets- og statsvidenskabelige: \emph{F. T. J. Gram};
i det lægevidenskabelige: \emph{C. Otto}; i det philosophiske: \emph{J. L.
Ussing}, og i det mathematisk-naturvidenskabelige: \emph{F. M. Liebmann}.

Festen aabnes og sluttes med den for denne Høitid bestemte Festsang, som udføres
af Studerende.
% p. 47
Videnskabernes og Universitetets Venner og Velyndere indbydes til, i Forening med
Universitetets Lærere og Studerende, ved deres Nærværelse at bidrage til Festens
Forskjønnelse.

\bigskip
\noindent Kjøbenhavn, den 31te October 1855.

\begin{center}---\end{center}

\begin{center}\textbf{Under Universitetets Segl.}\end{center}

% [University seal engraving: a crowned king enthroned, holding sword and sceptre,
%  above the Danish three-lions coat of arms, flanked by foliage]

\end{document}
